\documentclass[a4paper,12pt,twoside]{memoir}

% Castellano
\usepackage[spanish,es-tabla]{babel}
\selectlanguage{spanish}
\usepackage[utf8]{inputenc}
\usepackage[T1]{fontenc}
\usepackage{lmodern} % scalable font
\usepackage{microtype}
\usepackage{placeins}

\RequirePackage{booktabs}
\RequirePackage[table]{xcolor}
\RequirePackage{xtab}
\RequirePackage{multirow}

% Links
\PassOptionsToPackage{hyphens}{url}\usepackage[colorlinks]{hyperref}
\hypersetup{
    allcolors = {darkblue}
}
\usepackage{xcolor}
\definecolor{darkblue}{RGB}{0, 38, 73}

% Ecuaciones
\usepackage{amsmath}

% Rutas de fichero / paquete
\newcommand{\ruta}[1]{{\sffamily #1}}

% Párrafos
\nonzeroparskip

% Huérfanas y viudas
\widowpenalty100000
\clubpenalty100000

% Evitar solapes en el header
\nouppercaseheads

% Imagenes
\usepackage{graphicx}
\newcommand{\imagen}[3]{
	\begin{figure}[!h]
		\centering
		\includegraphics[width=#3\textwidth]{#1}
		\caption{#2}\label{fig:#1}
	\end{figure}
	\FloatBarrier
}

\newcommand{\imagenflotante}[2]{
	\begin{figure}%[!h]
		\centering
		\includegraphics[width=0.9\textwidth]{#1}
		\caption{#2}\label{fig:#1}
	\end{figure}
}



% El comando \figura nos permite insertar figuras comodamente, y utilizando
% siempre el mismo formato. Los parametros son:
% 1 -> Porcentaje del ancho de página que ocupará la figura (de 0 a 1)
% 2 --> Fichero de la imagen
% 3 --> Texto a pie de imagen
% 4 --> Etiqueta (label) para referencias
% 5 --> Opciones que queramos pasarle al \includegraphics
% 6 --> Opciones de posicionamiento a pasarle a \begin{figure}
\newcommand{\figuraConPosicion}[6]{%
  \setlength{\anchoFloat}{#1\textwidth}%
  \addtolength{\anchoFloat}{-4\fboxsep}%
  \setlength{\anchoFigura}{\anchoFloat}%
  \begin{figure}[#6]
    \begin{center}%
      \Ovalbox{%
        \begin{minipage}{\anchoFloat}%
          \begin{center}%
            \includegraphics[width=\anchoFigura,#5]{#2}%
            \caption{#3}%
            \label{#4}%
          \end{center}%
        \end{minipage}
      }%
    \end{center}%
  \end{figure}%
}

%
% Comando para incluir imágenes en formato apaisado (sin marco).
\newcommand{\figuraApaisadaSinMarco}[5]{%
  \begin{figure}%
    \begin{center}%
    \includegraphics[angle=90,height=#1\textheight,#5]{#2}%
    \caption{#3}%
    \label{#4}%
    \end{center}%
  \end{figure}%
}
% Para las tablas
\newcommand{\otoprule}{\midrule [\heavyrulewidth]}
%
% Nuevo comando para tablas pequeñas (menos de una página).
\newcommand{\tablaSmall}[5]{%
 \begin{table}
  \begin{center}
   \rowcolors {2}{gray!35}{}
   \begin{tabular}{#2}
    \toprule
    #4
    \otoprule
    #5
    \bottomrule
   \end{tabular}
   \caption{#1}
   \label{tabla:#3}
  \end{center}
 \end{table}
}

%
%Para el float H de tablaSmallSinColores
\usepackage{float}

%
% Nuevo comando para tablas pequeñas (menos de una página).
\newcommand{\tablaSmallSinColores}[5]{%
 \begin{table}[H]
  \begin{center}
   \begin{tabular}{#2}
    \toprule
    #4
    \otoprule
    #5
    \bottomrule
   \end{tabular}
   \caption{#1}
   \label{tabla:#3}
  \end{center}
 \end{table}
}

\newcommand{\tablaApaisadaSmall}[5]{%
\begin{landscape}
  \begin{table}
   \begin{center}
    \rowcolors {2}{gray!35}{}
    \begin{tabular}{#2}
     \toprule
     #4
     \otoprule
     #5
     \bottomrule
    \end{tabular}
    \caption{#1}
    \label{tabla:#3}
   \end{center}
  \end{table}
\end{landscape}
}

%
% Nuevo comando para tablas grandes con cabecera y filas alternas coloreadas en gris.
\newcommand{\tabla}[6]{%
  \begin{center}
    \tablefirsthead{
      \toprule
      #5
      \otoprule
    }
    \tablehead{
      \multicolumn{#3}{l}{\small\sl continúa desde la página anterior}\\
      \toprule
      #5
      \otoprule
    }
    \tabletail{
      \hline
      \multicolumn{#3}{r}{\small\sl continúa en la página siguiente}\\
    }
    \tablelasttail{
      \hline
    }
    \bottomcaption{#1}
    \rowcolors {2}{gray!35}{}
    \begin{xtabular}{#2}
      #6
      \bottomrule
    \end{xtabular}
    \label{tabla:#4}
  \end{center}
}

%
% Nuevo comando para tablas grandes con cabecera.
\newcommand{\tablaSinColores}[6]{%
  \begin{center}
    \tablefirsthead{
      \toprule
      #5
      \otoprule
    }
    \tablehead{
      \multicolumn{#3}{l}{\small\sl continúa desde la página anterior}\\
      \toprule
      #5
      \otoprule
    }
    \tabletail{
      \hline
      \multicolumn{#3}{r}{\small\sl continúa en la página siguiente}\\
    }
    \tablelasttail{
      \hline
    }
    \bottomcaption{#1}
    \begin{xtabular}{#2}
      #6
      \bottomrule
    \end{xtabular}
    \label{tabla:#4}
  \end{center}
}

%
% Nuevo comando para tablas grandes sin cabecera.
\newcommand{\tablaSinCabecera}[5]{%
  \begin{center}
    \tablefirsthead{
      \toprule
    }
    \tablehead{
      \multicolumn{#3}{l}{\small\sl continúa desde la página anterior}\\
      \hline
    }
    \tabletail{
      \hline
      \multicolumn{#3}{r}{\small\sl continúa en la página siguiente}\\
    }
    \tablelasttail{
      \hline
    }
    \bottomcaption{#1}
  \begin{xtabular}{#2}
    #5
   \bottomrule
  \end{xtabular}
  \label{tabla:#4}
  \end{center}
}



\definecolor{cgoLight}{HTML}{EEEEEE}
\definecolor{cgoExtralight}{HTML}{FFFFFF}

%
% Nuevo comando para tablas grandes sin cabecera.
\newcommand{\tablaSinCabeceraConBandas}[5]{%
  \begin{center}
    \tablefirsthead{
      \toprule
    }
    \tablehead{
      \multicolumn{#3}{l}{\small\sl continúa desde la página anterior}\\
      \hline
    }
    \tabletail{
      \hline
      \multicolumn{#3}{r}{\small\sl continúa en la página siguiente}\\
    }
    \tablelasttail{
      \hline
    }
    \bottomcaption{#1}
    \rowcolors[]{1}{cgoExtralight}{cgoLight}

  \begin{xtabular}{#2}
    #5
   \bottomrule
  \end{xtabular}
  \label{tabla:#4}
  \end{center}
}




\graphicspath{ {./img/} }

% Capítulos
\chapterstyle{bianchi}
\newcommand{\capitulo}[2]{
	\setcounter{chapter}{#1}
	\setcounter{section}{0}
	\setcounter{figure}{0}
	\setcounter{table}{0}
	\chapter*{#2}
	\addcontentsline{toc}{chapter}{#2}
	\markboth{#2}{#2}
}

% Apéndices
\renewcommand{\appendixname}{Apéndice}
\renewcommand*\cftappendixname{\appendixname}

\newcommand{\apendice}[1]{
	%\renewcommand{\thechapter}{A}
	\chapter{#1}
}

\renewcommand*\cftappendixname{\appendixname\ }

% Formato de portada
\makeatletter
\usepackage{xcolor}
\newcommand{\tutor}[1]{\def\@tutor{#1}}
\newcommand{\course}[1]{\def\@course{#1}}
\definecolor{cpardoBox}{HTML}{E6E6FF}
\def\maketitle{
  \null
  \thispagestyle{empty}
  % Cabecera ----------------
\noindent\includegraphics[width=\textwidth]{cabecera}\vspace{1cm}%
  \vfill
  % Título proyecto y escudo informática ----------------
  \colorbox{cpardoBox}{%
    \begin{minipage}{.8\textwidth}
      \vspace{.5cm}\Large
      \begin{center}
      \textbf{TFG del Grado en Ingeniería Informática}\vspace{.6cm}\\
      \textbf{\LARGE\@title{}}
      \end{center}
      \vspace{.2cm}
    \end{minipage}

  }%
  \hfill\begin{minipage}{.20\textwidth}
    \includegraphics[width=\textwidth]{escudoInfor}
  \end{minipage}
  \vfill
  % Datos de alumno, curso y tutores ------------------
  \begin{center}%
  {%
    \noindent\LARGE
    Presentado por \@author{}\\ 
    en Universidad de Burgos --- \@date{}\\
    Tutor: \@tutor{}\\
  }%
  \end{center}%
  \null
  \cleardoublepage
  }
\makeatother


% Datos de portada
\title{StatsGlowMind \\Documentación Técnica}
\author{Roberto Lastras Santos}
\tutor{José Manuel Galán Ordax}
\date{\today}

\begin{document}

\maketitle



\cleardoublepage



%%%%%%%%%%%%%%%%%%%%%%%%%%%%%%%%%%%%%%%%%%%%%%%%%%%%%%%%%%%%%%%%%%%%%%%%%%%%%%%%%%%%%%%%



\frontmatter


\clearpage

% Indices
\tableofcontents

\clearpage

\listoffigures

\clearpage

\listoftables

\clearpage

\mainmatter

\appendix

\apendice{Plan de Proyecto Software}

\section{Introducción}
En este apéndice de los anexos, plan de proyecto software, se ha realizado la planificación del proyecto, haciendo por una parte la planificación temporal y por otra el estudio de viabilidad económica y legal.

El objetivo de este anexo es ofrecer una visión general detallada de la gestión de nuestro proyecto software, de esta manera, este documento brinda como se ha planificado, ejecutado y controlado el proyecto desde el principio hasta el final.

\section{Planificación temporal}
Para la planificación del proyecto, se decidió utilizar una metodología ágil, por ello, tanto mi tutor como yo, hemos adoptado la metodología SCRUM para la gestión del proyecto, dividiendo el trabajo en sprints. 

Cada sprint tiene una duración de aproximadamente 2 semanas, lo que nos permite mantener un enfoque centrado en resultados y una capacidad de respuesta rápida a los cambios. Al finalizar el sprint, se ha realizado una tutoría para ver los avances en las tareas marcadas y proponer nuevas tareas para el siguiente sprint.

Para realizar scrum se ha utilizado Trello \cite{trello}, un software gratuito de gestión de tareas diseñado para facilitar la colaboración y el seguimiento de tareas de manera visual y colectiva.

En cuanto a la estimación de tiempo que se va a necesitar cada tarea he decidido usar esta tabla de relación entre las story point (puntos de historia) y tiempo en horas:

\begin{table}[!ht]
    \centering
    \begin{tabular}{|>{\centering\arraybackslash}m{4cm}|>{\centering\arraybackslash}m{4cm}|}
        \hline
        \rowcolor[rgb]{0.81,0.81,0.77}
        \textbf{Story Point} & \textbf{Estimación (horas)} \\
        \hline
        1 & 1 \\
        2 & 2 \\
        3 & 3.5 \\
        4 & 5 \\
        5 & 6.5 \\
        6 & 8 \\
        \hline
    \end{tabular}
    \caption{Estimaciones de tiempo para cada Story Point.}
    \label{tabla:estimaciones}
\end{table}

A continuación, se muestra cada sprint realizado a lo largo del proyecto, junto con una descripción sobre las tareas realizadas en el sprint, un gráfico de burndown para visualizar el progreso en la finalización de las tareas planificadas y una tabla que enumera todas las tareas planificadas con su estimación de story points.

\hfill

\subsection{Sprint 1 (24/01/2024 - 05/02/2024)}
En el primer sprint se han realizado los primeros pasos para realizar el TFG, entre ellos era decidir que metodología a utilizar e instalar y configurar el softaware que se ha usado en el proyecto, tanto en la gestión del proyecto: Trello, Zotero y GitHub; como en el desarrollo del proyecto: Angular y Firebase.

Además de comenzar la documentación y los estudios relacionado con el proyecto para así tener varias ideas y conclusiones antes de empezar a desarrollar la aplicación.

El gráfico de burndown del script y la tabla con las tareas y sus story points, se muestran a continuación:

\hfill

\imagen{anexos/burndown/burndown_sprint1}{Gráfico Burndown Sprint 1}{1}

\hfill

\begin{table}[!ht]
    \centering
    \resizebox{15cm}{!} {
    \begin{tabular}{|c|c|c|}
    \hline
    \rowcolor[rgb]{0.81,0.81,0.77}
    \textbf{Tareas} &\textbf{Etiquetas} & \textbf{Story points} \\ 
    \hline
        Leer el Reglamento del TFG & \cellcolor[HTML]{79E277}Investigación & 1\\
    \hline
        Crear un tablero en Trello & \cellcolor[HTML]{A50000}\textcolor{white}{Gestión} & 1\\ 
    \hline 
        Crear cuenta en Zotero & \cellcolor[HTML]{A50000}\textcolor{white}{Gestión} & 1\\ 
    \hline    
        Ver un tutorial de Zotero & \cellcolor[HTML]{79E277}Investigación & 2\\ 
    \hline  
        Crear un repositorio en GitHub & \cellcolor[HTML]{A50000}\textcolor{white}{Gestión} & 1\\ 
    \hline  
        Crear y configurar proyecto en Angular y Firebase & \begin{tabular}[c]{@{}c@{}c@{}}\cellcolor[HTML]{79E277}Investigación\\ \cellcolor[HTML]{A50000}\textcolor{white}{Gestión}\\ \cellcolor[HTML]{C382DF}Programación\end{tabular} & 4\\ 
    \hline
        Estudios relacionados con el TFG & \begin{tabular}[c]{@{}c@{}}\cellcolor[HTML]{79E277}Investigación\\ \cellcolor[HTML]{FC7E20}{Documentación}\end{tabular} & 3\\ 
    \hline
        Comenzar la documentación & \cellcolor[HTML]{FC7E20}{Documentación} & 2\\ 
    \hline
    \end{tabular}}
    \caption{Tareas Sprint 1}
    \label{tab:sprint1}
\end{table}

\clearpage

\subsection{Sprint 2 (06/02/2024 - 20/02/2024)}
En este sprint se han realizado las tareas relacionadas con los roles de usuario y la autenticación para los usuarios, para ello se ha implementado la herramienta de Firebase, Authentication, la cual ofrece un sistema de autenticación de usuarios seguro y fácil de usar.

Para implementar Authentication de Firebase en el proyecto, se ha creado el servicio Auth que permite registrar e iniciar sesión mediante correo electrónico o directamente mediante Google a un usuario. A parte se han creado los componentes necesarios para que un usuario pueda realizar el login, pueda registrarse y pueda acceder al main de la aplicación.

También se han realizado tareas de diseño como la realización del logotipo de la aplicación y la instalación de Bootstrap en el proyecto.

\hfill

\imagen{anexos/burndown/burndown_sprint2}{Gráfico Burndown Sprint 2}{1}

\clearpage

\begin{table}[!ht]
    \centering
    \resizebox{15cm}{!} {
    \begin{tabular}{|c|c|c|}
    \hline
    \rowcolor[rgb]{0.81,0.81,0.77}
    \textbf{Tareas} &\textbf{Etiquetas} & \textbf{Story points} \\ 
    \hline
        Agregar Autentificación en Firebase & \cellcolor[HTML]{A50000}\textcolor{white}{Gestión} & 2\\
    \hline
        Crear Servicio Auth (Authentication) & \cellcolor[HTML]{C382DF}Programación & 4\\ 
    \hline 
        Modificar AppComponent & \cellcolor[HTML]{C382DF}Programación & 2\\ 
    \hline    
        Crear Componente Register & \cellcolor[HTML]{C382DF}Programación & 4\\ 
    \hline  
        Crear Componente Login & \cellcolor[HTML]{C382DF}Programación & 3\\ 
    \hline  
        Crear Componente Main & \cellcolor[HTML]{C382DF}Programación & 2\\ 
    \hline
        Realizar diseño logotipo & \begin{tabular}[c]{@{}c@{}}\cellcolor[HTML]{6BBAFF}Diseño\\ \cellcolor[HTML]{C382DF}Programación\end{tabular} & 4\\ 
    \hline
        Instalar Bootstrap y Añadir CSS & \begin{tabular}[c]{@{}c@{}}\cellcolor[HTML]{6BBAFF}Diseño\\ \cellcolor[HTML]{C382DF}Programación\end{tabular} & 4\\ 
    \hline
    \end{tabular}}
    \caption{Tareas Sprint 2}
    \label{tab:sprint2}
\end{table}

\hfill

\subsection{Sprint 3 (21/02/2024 - 07/03/2024)}
El tercer sprint comienza las tareas que tienen que ver con las llamadas a las APIs de la NBA para mostrar información en la aplicación y también se realiza modificaciones y mejoras en cuanto a los roles de usuario y autenticación del sprint anterior.

Para conectar con las APIs, se ha creado un servicio llamado StatsService en el que se realizan las peticiones a las APIs, además se han creado los componentes Standings y Leaders con los que llaman al servicio Stats para recoger la información de las APIs y mostrarla.

En cuanto a las tareas de mejora de la autenticación de usuarios, se ha añadido a los usuarios un rol (usuario o administrador), se almacenan en Firestore (una base de datos NoSQL que proporciona Firebase) los usuarios que se han registrado y se comprueba si el usuario al hacer login, está registrado en Firestore.

\clearpage

\imagen{anexos/burndown/burndown_sprint3}{Gráfico Burndown Sprint 3}{1}


\begin{table}[!ht]
    \centering
    \resizebox{15cm}{!} {
    \begin{tabular}{|c|c|c|}
    \hline
    \rowcolor[rgb]{0.81,0.81,0.77}
    \textbf{Tareas} &\textbf{Etiquetas} & \textbf{Story points} \\ 
    \hline
        Añadir Rol de usuario en Register & \cellcolor[HTML]{C382DF}Programación & 2\\
    \hline
        Almacenar Usuarios en Firestore & \begin{tabular}[c]{@{}c@{}}\cellcolor[HTML]{79E277}Investigación\\ \cellcolor[HTML]{C382DF}Programación\end{tabular} & 3\\
    \hline
        Gestión de usuarios & \cellcolor[HTML]{C382DF}Programación & 4\\ 
    \hline 
        Investigar sobre peticiones a la API & \cellcolor[HTML]{79E277}Investigación & 3\\ 
    \hline    
        Conectar con la API y hacer routing & \cellcolor[HTML]{C382DF}Programación & 6\\ 
    \hline  
        Crear Componente Standings & \cellcolor[HTML]{C382DF}Programación & 3\\ 
    \hline  
        Crear Componente Leaders & \cellcolor[HTML]{C382DF}Programación & 2\\ 
    \hline  
        Comprobar si un usuario esta registrado en Firestore & \cellcolor[HTML]{C382DF}Programación & 2\\ 
    \hline
    \end{tabular}}
    \caption{Tareas Sprint 3}
    \label{tab:sprint3}
\end{table}

\hfill

\subsection{Sprint 4 (08/03/2024 - 21/03/2024)}
En este sprint 4 se instala e implementa Angular Material y se integran los componentes deseados en cada uno de los componentes del proyecto y se modifican todos los estilos para tener un diseño atractivo y responsive.

\hfill

\imagen{anexos/burndown/burndown_sprint4}{Gráfico Burndown Sprint 4}{1}

\hfill

\begin{table}[!ht]
    \centering
    \resizebox{15cm}{!} {
    \begin{tabular}{|c|c|c|}
    \hline
    \rowcolor[rgb]{0.81,0.81,0.77}
    \textbf{Tareas} &\textbf{Etiquetas} & \textbf{Story points} \\ 
    \hline
        Investigación Angular Material & \cellcolor[HTML]{79E277}Investigación & 3\\
    \hline
       Implementar Angular Material y personalizar Thema & \begin{tabular}[c]{@{}c@{}}\cellcolor[HTML]{6BBAFF}Diseño\\ \cellcolor[HTML]{C382DF}Programación\end{tabular} & 4\\
    \hline
        Integrar Material en MainComponent & \begin{tabular}[c]{@{}c@{}}\cellcolor[HTML]{6BBAFF}Diseño\\ \cellcolor[HTML]{C382DF}Programación\end{tabular} & 2\\ 
    \hline 
        Integrar Material en LeadersComponent & \begin{tabular}[c]{@{}c@{}}\cellcolor[HTML]{6BBAFF}Diseño\\ \cellcolor[HTML]{C382DF}Programación\end{tabular} & 3\\ 
    \hline    
       Integrar Material en LoginComponent & \begin{tabular}[c]{@{}c@{}}\cellcolor[HTML]{6BBAFF}Diseño\\ \cellcolor[HTML]{C382DF}Programación\end{tabular} & 2\\ 
    \hline  
        Integrar Material en AppComponent & \begin{tabular}[c]{@{}c@{}}\cellcolor[HTML]{6BBAFF}Diseño\\ \cellcolor[HTML]{C382DF}Programación\end{tabular} & 2\\
    \hline  
        Integrar Material en AdminComponent & \begin{tabular}[c]{@{}c@{}}\cellcolor[HTML]{6BBAFF}Diseño\\ \cellcolor[HTML]{C382DF}Programación\end{tabular} & 3\\
    \hline  
       Diseño Responsive & \begin{tabular}[c]{@{}c@{}}\cellcolor[HTML]{6BBAFF}Diseño\\ \cellcolor[HTML]{C382DF}Programación\end{tabular} & 4\\ 
    \hline
    \end{tabular}}
    \caption{Tareas Sprint 4}
    \label{tab:sprint4}
\end{table}

\clearpage

\subsection{Sprint 5 (22/03/2024 - 04/04/2024)}
En este sprint se ha dedicado en realizar apartados de la documentación de la memoria del proyecto y en conectar con las APIs para mostrar la información de los Partidos y de los Playoffs. Para ello se han creado esos componentes y se les ha agregado Angular Material y un estilo responsive.

\imagen{anexos/burndown/burndown_sprint5}{Gráfico Burndown Sprint 5}{0.9}

\begin{table}[!ht]
    \centering
    \resizebox{15cm}{!} {
    \begin{tabular}{|c|c|c|}
    \hline
    \rowcolor[rgb]{0.81,0.81,0.77}
    \textbf{Tareas} &\textbf{Etiquetas} & \textbf{Story points} \\ 
    \hline
        Documentación general de memoria &  \cellcolor[HTML]{FC7E20}{Documentación} & 3\\
    \hline
        Introducción del proyecto & \cellcolor[HTML]{FC7E20}{Documentación} & 2\\
    \hline
        Objetivos del proyecto & \cellcolor[HTML]{FC7E20}{Documentación} & 2\\
    \hline
        Estudios relacionados del proyecto & \cellcolor[HTML]{FC7E20}{Documentación} & 2\\
    \hline  
        Crear Componente Games & \cellcolor[HTML]{C382DF}Programación & 3\\ 
    \hline  
        Recoger datos del partido según la fecha & \cellcolor[HTML]{C382DF}Programación & 2\\
    \hline  
        Material y Responsive en GamesComponent & \begin{tabular}[c]{@{}c@{}}\cellcolor[HTML]{6BBAFF}Diseño\\ \cellcolor[HTML]{C382DF}Programación\end{tabular} & 3\\
    \hline  
        Crear Componente Playoffs & \cellcolor[HTML]{C382DF}Programación & 3\\ 
    \hline  
        Material y Responsive en PlayoffsComponent & \begin{tabular}[c]{@{}c@{}}\cellcolor[HTML]{6BBAFF}Diseño\\ \cellcolor[HTML]{C382DF}Programación\end{tabular} & 3\\
    \hline  
        Añadir sistema de filtrado en LeadersComponent & \begin{tabular}[c]{@{}c@{}}\cellcolor[HTML]{6BBAFF}Diseño\\ \cellcolor[HTML]{C382DF}Programación\end{tabular} & 3\\
    \hline
    \end{tabular}}
    \caption{Tareas Sprint 5}
    \label{tab:sprint5}
\end{table}

\clearpage

\subsection{Sprint 6 (05/04/2024 - 18/04/2024)}

\hfill

En el sprint 6, se han creado en la base de datos de Firestore las colecciones: Teams, Games y Players. 

En las cuales lo que he hecho ha sido primero realizar las peticiones adecuadas en la API, obtener los datos que nos devuelve la API y almacenarlos en la base de datos. Una vez estén todos los datos almacenados, se han realizado consultas para mostrar la información deseada en tablas u otros componentes, en la aplicación.

También se ha dedicado gran parte del sprint a la realización de la documentación de la memoria (Conceptos teóricos y Técnicas y Herramientas) y la documentación del anexo Plan de Proyecto Software (Introducción y Planificación temporal).

\hfill

\imagen{anexos/burndown/burndown_sprint6}{Gráfico Burndown Sprint 6}{1}

\clearpage

\begin{table}[!ht]
    \centering
    \resizebox{15cm}{!} {
    \begin{tabular}{|c|c|c|}
    \hline
    \rowcolor[rgb]{0.81,0.81,0.77}
    \textbf{Tareas} &\textbf{Etiquetas} & \textbf{Story points} \\ 
    \hline
        Creación de colecciones Firestore &  \begin{tabular}[c]{@{}c@{}}\cellcolor[HTML]{A50000}\textcolor{white}{Gestión}\\ \cellcolor[HTML]{C382DF}Programación\end{tabular} & 1\\
    \hline
        Almacenar en Firestore los equipos & \cellcolor[HTML]{C382DF}Programación & 4\\
    \hline
        Material y Responsive en TeamsComponent & \begin{tabular}[c]{@{}c@{}}\cellcolor[HTML]{6BBAFF}Diseño\\ \cellcolor[HTML]{C382DF}Programación\end{tabular} & 4\\
    \hline
        Almacenar en Firestore las estadísticas de los equipos & \cellcolor[HTML]{C382DF}Programación & 6\\
    \hline  
        Estilos y Responsive en TeamStatsComponent & \begin{tabular}[c]{@{}c@{}}\cellcolor[HTML]{6BBAFF}Diseño\\ \cellcolor[HTML]{C382DF}Programación\end{tabular} & 2\\ 
    \hline  
        Almacenar en Firestore los jugadores & \cellcolor[HTML]{C382DF}Programación & 3\\
    \hline  
        Estilos y Responsive en TeamPlayersComponent & \begin{tabular}[c]{@{}c@{}}\cellcolor[HTML]{6BBAFF}Diseño\\ \cellcolor[HTML]{C382DF}Programación\end{tabular} & 2\\
    \hline  
        Documentación Memoria & \cellcolor[HTML]{FC7E20}{Documentación} & 6\\ 
    \hline  
        Documentación Anexos & \cellcolor[HTML]{FC7E20}{Documentación} & 6\\
    \hline  
        Creación de un Footer & \begin{tabular}[c]{@{}c@{}}\cellcolor[HTML]{6BBAFF}Diseño\\ \cellcolor[HTML]{C382DF}Programación\end{tabular} & 3\\
    \hline
    \end{tabular}}
    \caption{Tareas Sprint 6}
    \label{tab:sprint6}
\end{table}

\hfill


\subsection{Sprint 7 (19/04/2024 - 06/05/2024)}
En este sprint se ha enfocado principalmente en el análisis de los jugadores, de manera que se han creado el componente AnalysisComponent, donde se muestra una tabla con todos los jugadores de la NBA, junto con un buscador para que sea más sencillo filtrar los jugadores y además unas Cards con los jugadores más relevantes de la temporada; y el componente ChartsPlayerComponent donde se muestra las estadísticas media de la temporada del jugador y un conjunto de gráficos que muestran y analizan las estadísticas del jugador durante la liga. Para mostrar los gráficos se ha implementado la librería Chart.js que nos permite insertar gráficos de varios tipos de manera fácil y visual.

A parte del análisis de los jugadores, también se ha realizado un nuevo diseño del main para que se visualizara más estético e intuitivo para el usuario. Además de avanzar con la documentación de la memoria y de los anexos del proyecto.

\hfill

\imagen{anexos/burndown/burndown_sprint7}{Gráfico Burndown Sprint 7}{1}

\hfill

\begin{table}[!ht]
    \centering
    \resizebox{15cm}{!} {
    \begin{tabular}{|c|c|c|}
    \hline
    \rowcolor[rgb]{0.81,0.81,0.77}
    \textbf{Tareas} &\textbf{Etiquetas} & \textbf{Story points} \\ 
    \hline
        Creación del componente Analysis &  \cellcolor[HTML]{C382DF}Programación & 6\\
    \hline
        Estilos y Responsive en AnalysisComponent & \begin{tabular}[c]{@{}c@{}}\cellcolor[HTML]{6BBAFF}Diseño\\ \cellcolor[HTML]{C382DF}Programación\end{tabular} & 2\\
    \hline
        Creación del componente ChartsPlayer &  \cellcolor[HTML]{C382DF}Programación & 5\\
    \hline
        Almacenar en Firestore las estadísticas de cada partido & \cellcolor[HTML]{C382DF}Programación & 5\\
    \hline  
        Calcular estadísticas medias del jugador & \cellcolor[HTML]{C382DF}Programación & 3\\
    \hline  
        Implementar Chart.js & \begin{tabular}[c]{@{}c@{}}\cellcolor[HTML]{6BBAFF}Diseño\\ \cellcolor[HTML]{C382DF}Programación\\ \cellcolor[HTML]{79E277}Investigación\end{tabular} & 3\\ 
    \hline  
        Estilos y Responsive en ChartsPlayerComponent & \begin{tabular}[c]{@{}c@{}}\cellcolor[HTML]{6BBAFF}Diseño\\ \cellcolor[HTML]{C382DF}Programación \end{tabular} & 2\\
    \hline  
        Rediseñar MainComponent & \begin{tabular}[c]{@{}c@{}}\cellcolor[HTML]{6BBAFF}Diseño\\ \cellcolor[HTML]{C382DF}Programación\end{tabular} & 6\\
    \hline  
        Documentación Memoria & \cellcolor[HTML]{FC7E20}{Documentación} & 6\\ 
    \hline  
        Documentación Anexos & \cellcolor[HTML]{FC7E20}{Documentación} & 6\\
    \hline
    \end{tabular}}
    \caption{Tareas Sprint 7}
    \label{tab:sprint7}
\end{table}

\clearpage

\subsection{Sprint 8 (07/05/2024 - 22/05/2024)}
El sprint 8 ha sido uno de los más complicados ya que se ha realizado una investigación exhaustiva sobre como implementar modelos de machine learning en mi proyecto, después de investigar durante bastante tiempo, decidí que la mejor opción era utilizar la funcionalidad Cloud Functions de Firebase, que permite subir funciones en la nube sin necesidad de crear un servidor backend en mi proyecto.

Una vez realizada la investigación, me dedique a entrenar mediante los datos del histórico de las estadísticas de cada partido de la temporada de la NBA, un modelo utilizando el algoritmo Random Forest. Después cree la función que realiza la predicción de los partidos utilizando este modelo y mostré la predicción en mi componente mediante un gráfico de Chart.js.

\hfill

\imagen{anexos/burndown/burndown_sprint8}{Gráfico Burndown Sprint 8}{1}

\hfill

\begin{table}[!ht]
    \centering
    \resizebox{15cm}{!} {
    \begin{tabular}{|c|c|c|}
    \hline
    \rowcolor[rgb]{0.81,0.81,0.77}
    \textbf{Tareas} & \textbf{Etiquetas} & \textbf{Story points} \\ 
    \hline
        Investigación sobre modelos de ML &  \cellcolor[HTML]{79E277}Investigación & 6\\
    \hline
        Generar tabla de datos para el modelo & \begin{tabular}[c]{@{}c@{}}\cellcolor[HTML]{A50000}\textcolor{white}{Gestión}\\ \cellcolor[HTML]{C382DF}Programación\end{tabular} & 3\\
    \hline
        Entrenar modelo Random Forest &  \cellcolor[HTML]{C382DF}Programación & 4\\
    \hline
        Implementar Firebase Cloud Functions & \begin{tabular}[c]{@{}c@{}}\cellcolor[HTML]{A50000}\textcolor{white}{Gestión}\\ \cellcolor[HTML]{C382DF}Programación\end{tabular} & 2\\
    \hline  
        Crear la función de predecir en Python & \begin{tabular}[c]{@{}c@{}}\cellcolor[HTML]{79E277}Investigación\\ \cellcolor[HTML]{C382DF}Programación\end{tabular} & 6\\
    \hline  
        Crear servicio PredictionService & \cellcolor[HTML]{C382DF}Programación & 2\\ 
    \hline  
        Crear componente PredictComponent & \cellcolor[HTML]{C382DF}Programación & 3\\
    \hline  
        Mostrar predicción con Chart.js & \begin{tabular}[c]{@{}c@{}}\cellcolor[HTML]{6BBAFF}Diseño\\ \cellcolor[HTML]{C382DF}Programación\end{tabular} & 3\\
    \hline  
        Estilos y responsive en PredictComponent & \begin{tabular}[c]{@{}c@{}}\cellcolor[HTML]{6BBAFF}Diseño\\ \cellcolor[HTML]{C382DF}Programación\end{tabular} & 4\\
    \hline  
        Documentación Memoria & \cellcolor[HTML]{FC7E20}{Documentación} & 6\\ 
    \hline  
        Documentación Anexos & \cellcolor[HTML]{FC7E20}{Documentación} & 6\\
    \hline
    \end{tabular}}
    \caption{Tareas Sprint 8}
    \label{tab:sprint8}
\end{table}

\hfill

\subsection{Sprint 9 (23/05/2024 - 07/06/2024)}

Este último sprint, se dedicó principalmente a internacionalizar el proyecto, pudiendo así cambiar de idioma de español a inglés y viceversa. También se implementaron algunas mejoras en la aplicación, como la creación de componentes nuevos para proporcionar más información acerca de la aplicación a los usuarios.

En el último tramo del sprint tan solo se terminó de documentar la memoria y los anexos del proyecto, esto llevo bastante tiempo ya que es mucha información que documentar.

\imagen{anexos/burndown/burndown_sprint9}{Gráfico Burndown Sprint 9}{1}

\begin{table}[!ht!]
    \centering
    \resizebox{15cm}{!} {
    \begin{tabular}{|c|c|c|}
    \hline
    \rowcolor[rgb]{0.81,0.81,0.77}
    \textbf{Tareas} &\textbf{Etiquetas} & \textbf{Story points} \\ 
    \hline
        Crear componente AboutComponent &  \begin{tabular}[c]{@{}c@{}}\cellcolor[HTML]{6BBAFF}Diseño\\ \cellcolor[HTML]{C382DF}Programación\end{tabular} & 3\\
    \hline
        Crear componente GlossaryComponent &  \begin{tabular}[c]{@{}c@{}}\cellcolor[HTML]{6BBAFF}Diseño\\ \cellcolor[HTML]{C382DF}Programación\end{tabular} & 3\\
    \hline
        Mejorar el Footer &  \begin{tabular}[c]{@{}c@{}}\cellcolor[HTML]{6BBAFF}Diseño\\ \cellcolor[HTML]{C382DF}Programación\end{tabular} & 1\\
    \hline
        Añadir select para cambiar el idioma & \begin{tabular}[c]{@{}c@{}}\cellcolor[HTML]{6BBAFF}Diseño\\ \cellcolor[HTML]{C382DF}Programación\end{tabular} & 2\\
    \hline  
        Modificar usuarios siendo admin & \cellcolor[HTML]{C382DF}Programación & 2\\ 
    \hline  
        Investigación sobre internacionalización & \cellcolor[HTML]{79E277}Investigación & 6\\
    \hline  
        Internacionalizar el proyecto & \begin{tabular}[c]{@{}c@{}}\cellcolor[HTML]{A50000}\textcolor{white}{Gestión}\\ \cellcolor[HTML]{C382DF}Programación\end{tabular} & 6\\
    \hline  
        Documentación Memoria & \cellcolor[HTML]{FC7E20}{Documentación} & 6\\ 
    \hline  
        Documentación Anexos I & \cellcolor[HTML]{FC7E20}{Documentación} & 6\\
    \hline  
        Documentación Anexos II & \cellcolor[HTML]{FC7E20}{Documentación} & 6\\
    \hline  
        Documentación GitHub & \cellcolor[HTML]{FC7E20}{Documentación} & 4\\
    \hline
    \end{tabular}}
    \caption{Tareas Sprint 9}
    \label{tab:sprint9}
\end{table}

\clearpage

\section{Estudio de viabilidad}
En esta sección se realiza un estudio exhaustivo de viabilidad económica y legal para garantizar la viabilidad y sustentabilidad del proyecto a largo plazo.

\hfill

\subsection{Viabilidad económica}

A continuación, se detallan los costes y beneficios asociados con la aplicación en caso de que se lanzara al mercado.

\subsubsection{Costes}

\hfill

\textbf{Coste de empleados}

La aplicación ha sido desarrollada por un único empleado y se calcula que se han realizado unas 360 horas de trabajo repartidas durante 4 meses y medio (18 semanas). Esto supone una carga de trabajo de unas 20 horas semanales. El salario del alumno se estima que va a ser de unos 15\text{€}/hora, el cálculo del salario bruto laboral mensual sería el siguiente:

$$ 20\frac{horas}{semana}\times15\frac{\text{€}}{hora}\times4\frac{semanas}{mes}=1200\text{€}\hspace{0.5em}al\hspace{0.5em}mes  $$

Para calcular el salario real que recibirá el empleado se debe calcular los impuestos que paga la empresa por el empleado. Esto se puede consultar en la página oficial de la seguridad social \cite{ss}.

\hfill

Los impuestos son:
\begin{itemize}
\tightlist
    \item 23.60\% de contingencias.
    \item 5.50\% de desempleo.
    \item 0.20\% de FOGASA.
    \item 0.60\% de formación profesional.
\end{itemize}

\imagen{anexos/impuestos}{Régimen general de la Seguridad Social}{1}

Por lo tanto, teniendo en cuenta estos impuestos, calculamos que la empresa tiene que pagar por el empleado:

$$ \frac{1200\frac{\text{€}}{mes}}{1-(0.236+0.055+0.002+0.006)}=1711,84\text{€}\hspace{0.5em}al\hspace{0.5em}mes $$

Por otra parte, también se cuenta con un profesor tutelando al alumno, con amplios conocimientos en el tema, lo que supone un salario de unos 40€/hora. El tutor trabaja unas 2 horas por cada dos semanas, es decir, 1 hora a la semana. Por lo tanto:

$$ 1\frac{hora}{semana}\times40\frac{\text{€}}{hora}\times4\frac{semanas}{mes}=160\text{€}\hspace{0.5em}al\hspace{0.5em}mes $$

Este cálculo corresponde con el salario bruto del profesor, al cual hay que sumar los impuestos:

$$ \frac{160\frac{\text{€}}{mes}}{1-(0.236+0.055+0.002+0.006)}=228,25\text{€}\hspace{0.5em}al\hspace{0.5em}mes $$

Finalmente obtenemos que la empresa deberá pagar al mes 1940.10\text{€}. Y como el proyecto ha durado cuatro meses y medio, el coste total será de 8730,45\text{€}.

\hfill

\textbf{Coste de Hardware}

Los recursos hardware utilizados para el desarrollo del proyecto, ha sido únicamente un ordenador portátil cuyo coste fue de 950\text{€}. Sin embargo, después de todo el uso recibido durante los 4 años de la carrera, se puede decir que el portátil ya ha sido amortizado en su totalidad, por lo tanto, los costes de recursos hardware han sido de un total de 0\text{€}.

\hfill

\textbf{Coste de Software}

Para el proyecto todas las herramientas software utilizadas han sido gratuitas excepto sistema operativo (Windows 10 Pro) del ordenador portátil, el cual ya está amortizado. Por lo que el coste de software real ha sido de 0\text{€}.

Pero para un mejor rendimiento de la aplicación, lo ideal sería cambiar el plan gratuito de la API utilizada por otro plan de pago que permite realizar muchísimas más peticiones por \$25,00 al mes, que serían 23\text{€} al mes. Y también modificar el plan de facturación de Firebase, para poder realizar muchas más operaciones de lectura y escritura  en la base de datos, añadir más GB almacenados en el servicio hosting o incrementar las invocaciones a la función de Cloud Functions; esto supondría un gasto de unos \$1000,00 al mes que son 919,81\text{€} al mes. Por lo que el costo de software ideal sería de 942,81\text{€} al mes.

Ya que es preferible añadir estos costes de software para mejorar la aplicación, se va a calcular el coste total del proyecto a partir del tercer mes (este incluido), que es cuando se comienza a utilizar ambos softwares constantemente. Ambos softwares tienen una amortización de 4 años, y como empezamos a contar a partir del tercer mes, que serían 2 meses y medio, costaría un total de 2357,03\text{€}.

$$ \frac{2357,03\text{€}}{4\hspace{0.5em}años}=589,26\text{€}\hspace{0.5em}al\hspace{0.5em}año $$

Como este plan de software dura 2 meses y medio, se divide 589,26\text{€} entre 4,8, por lo que el coste total del software sería de 122,76\text{€}.

\hfill

\textbf{Coste Total}

Al calcular los costes totales del proyecto, también hay que tener en cuenta los costes indirectos, como es la tarifa de internet, de la luz, de la infraestructura, \ldots. Este valor se establece como el 15\% de los gastos totales. De esta forma obtenemos un gasto total del proyecto de:

\begin{table}[!ht]
    \centering
    \begin{tabular}{|>{\centering\arraybackslash}m{4cm}|>{\centering\arraybackslash}m{4cm}|}
        \hline
        \rowcolor[rgb]{0.81,0.81,0.77}
        \textbf{Costes} & \textbf{Total \text{€}} \\
        \hline
        Coste de empleados & 8730,45 \text{€} \\
        Coste de hardware & 0 \text{€} \\
        Coste de software & 122,76 \text{€} \\
        Coste indirecto & 1327,98 \text{€} \\
        \hline
        \textbf{Total} & 10181,19 \text{€} \\
        \hline
    \end{tabular}
    \caption{Costes totales.}
    \label{tabla:costes}
\end{table}

\subsubsection{Beneficios}
Esta aplicación se ha desarrollado con carácter educativo, por lo que no se obtiene ningún beneficio del uso de la aplicación.

\clearpage

\subsection{Viabilidad legal}

En este apartado se va a realizar el estudio de las licencias que tienen las librerías y herramientas utilizadas para desarrollar el proyecto, como la licencia del mismo. En este proyecto se ha utilizado las siguientes herramientas y librerías:

\begin{table}[!ht]
    \centering
    \begin{tabular}{|>{\centering\arraybackslash}m{7cm}|>{\centering\arraybackslash}m{2cm}|>{\centering\arraybackslash}m{3cm}|}
        \hline
        \rowcolor[rgb]{0.81,0.81,0.77}
        \textbf{Herramienta/Librería} & \textbf{Versión} & \textbf{Licencia}\\
        \hline
        Angular & 17.2.3 & MIT \\
        AngularCLI & 17.0.1 & MIT \\
        TypeScript & 5.4.3 & Apache 2 \\
        npm & 10.1.0 & Artistic 2 \\
        Firebase & 10.9.0 & Apache 2 \\
        rxjs & 7.8.0 & Apache 2 \\
        Bootstrap & 5.3.2 & MIT \\
        \text{@}angular/material & 17.3.1 & MIT \\
        Chart.js & 4.4.2 & MIT \\
        ngx-toastr & 18.0.0 & MIT \\
        PapaParse & 5.4.1 & MIT \\
        \text{@}ngx-translate/core & 15.0.0 & MIT \\
        firebase-admin & 6.5.0 & Apache 2 \\
        firebase-functions & 0.4.1 & Apache 2 \\
        google-cloud-storage & 2.16.0 & Apache 2 \\
        scikit-learn & 1.4.2 & BSD \\
        pandas & 2.2.2 & BSD \\
        Flask & 3.0.3 & BSD \\
        flask cors & 4.0.1 & MIT \\
        \hline
    \end{tabular}
    \caption{Licencias y Versiones.}
    \label{tabla:licencias}
\end{table}

Todas las herramientas y librerías utilizadas en el desarrollo del proyecto tienen licencias de software libre, por lo cual, no existe ninguna clase de restricción o limitación por la que se deba aplicar una licencia más restrictiva a este proyecto.

Por esta razón, se ha decidido aplicar la licencia MIT en el proyecto, de tal manera que se permite su uso de forma libre siempre y cuando se proporcione atribución al autor original y se incluya un aviso de derechos de autor en todas las copias del software.
\apendice{Especificación de Requisitos}

\section{Introducción}

En este apéndice, se mostrarán y desarrollarán los objetivos y requisitos que debe cumplir la aplicación según han sido establecidos al comienzo del proyecto y los que se han ido añadiendo durante el proyecto.


\section{Objetivos generales}
Al comienzo del proyecto se establecieron los siguientes objetivos generales:
\begin{itemize}
\tightlist
    \item 
        Desarrollar una la aplicación web utilizando Angular como framework y Firebase como plataforma de alojamiento para la base de datos, autenticación de usuarios y servicio de hosting.
    \item
        Implementar la integración con una API con información de la NBA para recopilar datos en tiempo real. 
    \item 
        Gestionar la base de datos NoSQL (Firestore) para almacenar y gestionar los datos de la aplicación web.
    \item 
        Implementar un sistema de autenticación de usuarios.
    \item 
        Proporcionar información y estadísticas detalladas sobre la actual temporada de la NBA. 
    \item 
        Realizar análisis de rendimiento de jugadores mediante estadísticas y gráficas de datos.
    \item 
        Generar predicciones de resultados de partidos a partir de estadísticas y resultados de cada partido.
\end{itemize}


\section{Catálogo de requisitos}
En esta sección se exponen los requisitos funcionales y no funcionales del proyecto.

\subsection{Requisitos funcionales}
\begin{itemize}
\tightlist
    \item \textbf{RF1 - Login de usuarios:} La aplicación debe ser capaz de permitir a un usuario autenticarse y poder acceder a ciertas funcionalidades de la aplicación.
    \item \textbf{RF2 - Registro de usuarios:} La aplicación debe ser capaz de permitir a un usuario registrarse y almacenar sus datos en la base de datos.
    \item \textbf{RF3 - Logout de usuarios:} La aplicación debe permitir cerrar sesión a un usuario.
    \item \textbf{RF4 - Gestión de usuarios:} Un usuario con rol de administrador, debe poder gestionar el resto de usuarios, así como eliminar y modificar sus cuentas.
    \item \textbf{RF5 - Menú de la aplicación:} La aplicación debe contener un menú en el que puedas navegar por las distintas funcionalidades que tiene la aplicación.
    \item \textbf{RF6 - Main de la aplicación:} La aplicación debe de tener una página principal (main) en la que se muestren las funcionalidades de la aplicación de manera visual e intuitiva.
    \item \textbf{RF7 - Footer de la aplicación:} La aplicación debe contener un footer en el que se muestre más información sobre la aplicación.
        \begin{itemize}
        \tightlist
            \item \textbf{RF7.1 - Desarrollo de la aplicación:} En el footer debe contener información sobre herramientas y técnicas que se han utilizado para desarrollar la aplicación.
            \item \textbf{RF7.2 - Diseño de la aplicación:} En el footer debe contener información sobre herramientas y técnicas que se han utilizado para diseñar la aplicación.
            \item \textbf{RF7.3 - Acerca de la aplicación:} En el footer debe contener información sobre la aplicación.
            \item \textbf{RF7.4 - Redes sociales:} En el footer debe contener información sobre las redes sociales de la aplicación.
        \end{itemize}
    
    \item \textbf{RF8 - Mostrar información en tiempo real:} La aplicación debe mostrar información y estadísticas detalladas sobre la actual temporada de la NBA, en tiempo real.
        \begin{itemize}
        \tightlist
            \item \textbf{RF8.1 - Mostrar clasificación:} La aplicación debe poder mostrar la clasificación de la NBA junto con sus estadísticas.
            \item \textbf{RF8.2 - Mostrar playoffs:} La aplicación debe poder mostrar la eliminatoria de los playoffs de la NBA.
            \item \textbf{RF8.3 - Mostrar partidos:} La aplicación debe poder mostrar los partidos que se van a disputar y los que se han disputado, con sus resultados.
            \item \textbf{RF8.4 - Mostrar líderes:} La aplicación debe poder mostrar una tabla con los líderes de la NBA junto con sus estadísticas, y poder clasificar la lista según la categoría de líderes deseada.
        \end{itemize}
    \item \textbf{RF9 - Visualización de equipos:} La aplicación debe ser capaz de mostrar todos los equipos de todas las divisiones de cada conferencia de la NBA.
        \begin{itemize}
        \tightlist
            \item \textbf{RF9.1 - Estadísticas del equipo:} La aplicación debe poder mostrar mediante tablas, todas las estadísticas disponibles del equipo seleccionado.
            \item \textbf{RF9.2 - Jugadores del equipo:} La aplicación debe poder mostrar mediante una tabla, todos los jugadores junto con su información, del equipo seleccionado.
        \end{itemize}
    \item \textbf{RF10 - Análisis de jugadores:}
        \begin{itemize}
        \tightlist
            \item \textbf{RF10.1 - Mostrar jugadores:} La aplicación debe poder mostrar mediante una tabla todos los jugadores de la NBA y mediante cards los jugadores más destacados de la temporada.
            \item \textbf{RF10.2 - Gráficas y estadísticas del jugador:} La aplicación debe poder mostrar mediante una tabla las estadísticas medias de la temporada y analizar estas estadísticas mediante gráficas.
        \end{itemize}
    \item \textbf{RF11 - Predicción de partidos:} La aplicación debe ser capaz de proporcionar una predicción válida de un partido entre dos equipos seleccionados por el usuario.
    \item \textbf{RF12 - Cambio de idioma:} La aplicación debe poder modificar la configuración del idioma, de español a inglés y viceversa, si el usuario lo desea.
\end{itemize}

\subsubsection{Requisitos no funcionales}
\begin{itemize}
\tightlist
    \item \textbf{RNF1 - Usabilidad:} La aplicación debe tener una interfaz visual e intuitiva
    \item \textbf{RNF2 - Responsividad:} La aplicación debe ser capaz de adaptarse a los distintos tamaños de pantalla.
    \item \textbf{RNF3 - Internacionalización:} La aplicación debe poder configurar el idioma establecido.
    \item \textbf{RNF4 - Escalabilidad:} La aplicación debe estar preparada para desarrollar nuevas funcionalidades sin aumentar la carga de trabajo.
    \item \textbf{RNF5 - Mantenibilidad:} La aplicación debe permitir realizar mantenimientos y actualizaciones de forma eficaz. 
    \item \textbf{RNF6 - Seguridad:} La aplicación debe de proteger la integridad y confidencialidad de los datos de los usuarios y ocultar las claves y credenciales de acceso.
    \item \textbf{RNF7 - Disponibilidad:} La aplicación debe estar disponible para los usuarios todo el tiempo posible.
    
\end{itemize}



\begin{table}[p]
\section{Especificación de requisitos}

\hfill

\subsection{Casos de uso}

\hfill

	\centering
	\begin{tabularx}{\linewidth}{ p{0.21\columnwidth} p{0.71\columnwidth}}
		\toprule
		\textbf{CU-1}    & \textbf{Login de usuarios}\\
		\toprule
		\textbf{Requisitos asociados} & RF1 \\
		\textbf{Descripción}          & Permite a un usuario iniciar sesión y poder acceder a todas las funcionalidades de la aplicación y en caso de ser un usuario con rol de Administrador, también tiene acceso a la gestión de usuarios. \\
		\textbf{Precondición}         & Tener cuenta de Google si el usuario desea iniciar sesión mediante Google.\\
		\textbf{Acciones}             &
		\begin{enumerate}
			\def\labelenumi{\arabic{enumi}.}
			\tightlist
			\item El usuario accede a la página "\textit{Iniciar Sesión}".
			\item Opción 1: El usuario rellena el formulario introduciendo su correo electrónico y contraseña.
			Pulsa el botón de "\textit{Enviar}".
			\item Opción 2: El usuario pulsa el botón "\textit{Login With Google}".
			Selecciona su cuenta de Google.
		\end{enumerate}\\
		\textbf{Postcondición}        & El registro del usuario deber estar almacenado en la base de datos.		\\
		\textbf{Excepciones}          & 
		\begin{enumerate}
			\def\labelenumi{\arabic{enumi}.}
			\tightlist
			\item El usuario no exista en la base de datos (\textit{mensaje}).
			\item La contraseña no sea válida (\textit{mensaje}).
			\item El correo electrónico no sea válido (\textit{mensaje}).		
            \end{enumerate}\\
		\textbf{Importancia}          & Alta \\
		\bottomrule
	\end{tabularx}
	\caption{CU-1 Login de usuarios.}
\end{table}

\begin{table}[p]
	\centering
	\begin{tabularx}{\linewidth}{ p{0.21\columnwidth} p{0.71\columnwidth} }
		\toprule
		\textbf{CU-2}    & \textbf{Registro de usuarios}\\
		\toprule
		\textbf{Requisitos asociados} & RF2 \\
		\textbf{Descripción}          & Permite a un usuario crearse una cuenta mediante Correo electrónico/contraseña o mediante la cuenta de Google. \\
		\textbf{Precondición}         & Tener cuenta de Google si el usuario desea registrarse mediante Google. \\
		\textbf{Acciones}             &
		\begin{enumerate}
			\def\labelenumi{\arabic{enumi}.}
			\tightlist
			\item El usuario accede a la página "\textit{Iniciar Sesión}".
			\item El usuario selecciona a la opción "\textit{Registrarse}".
			\item Opción 1: El usuario rellena el formulario introduciendo su nombre de usuario, correo electrónico y contraseña, repite la contraseña y selecciona el rol de usuario.
			Pulsa el botón de "\textit{Enviar}".
			\item Opción 2: El usuario pulsa el botón "\textit{Register With Google}".
			Selecciona su cuenta de Google.
		\end{enumerate}\\
		\textbf{Postcondición}        &  El registro del usuario tiene que estar autenticado en Firebase. \\
		\textbf{Excepciones}          & 
		\begin{enumerate}
			\def\labelenumi{\arabic{enumi}.}
			\tightlist
			\item El usuario ya exista en la base de datos (\textit{mensaje}).
                \item Las contraseñas no coincidan (\textit{mensaje}).
			\item La contraseña no sea válida (\textit{mensaje}).
			\item El correo electrónico no sea válido (\textit{mensaje}).
            \end{enumerate}\\
		\textbf{Importancia}          & Alta \\
		\bottomrule
	\end{tabularx}
	\caption{CU-2 Registro de usuarios.}
\end{table}

\begin{table}[p]
	\centering
	\begin{tabularx}{\linewidth}{ p{0.21\columnwidth} p{0.71\columnwidth} }
		\toprule
		\textbf{CU-3}    & \textbf{Logout de usuarios}\\
		\toprule
		\textbf{Requisitos asociados} & RF3 \\
		\textbf{Descripción}          & Permite a un usuario cerrar sesión de la cuenta que tenga activa. \\
		\textbf{Precondición}         & Estar con la sesión iniciada. \\
		\textbf{Acciones}             & El usuario pulsa el botón "\textit{Cerrar Sesión}". \\
		\textbf{Postcondición}        & -	\\
		\textbf{Excepciones}          & -	\\
		\textbf{Importancia}          & Baja \\
		\bottomrule
	\end{tabularx}
	\caption{CU-3 Logout de usuarios.}
\end{table}

\begin{table}[p]
	\centering
	\begin{tabularx}{\linewidth}{ p{0.21\columnwidth} p{0.71\columnwidth} }
		\toprule
		\textbf{CU-4}    & \textbf{Gestión de usuarios}\\
		\toprule
		\textbf{Requisitos asociados} & RF4 \\
		\textbf{Descripción}          & Permite a un usuario con rol de administrador, modificar datos de las cuentas de usuario y eliminar usuarios. \\
		\textbf{Precondición}         & El usuario debe estar registrado y tener permisos de administrador.\\
		\textbf{Acciones}             &
		\begin{enumerate}
			\def\labelenumi{\arabic{enumi}.}
			\tightlist
			\item El usuario administrador accede a la página pulsando su "\textit{Nombre de Usuario}".
			\item Si el usuario administrador desea modificar los datos de un usuario, selecciona a la opción "\textit{Modificar}", cambia los datos deseados y pulsa en "\textit{Confirmar}".
			\item Si el usuario administrador desea eliminar a un usuario, selecciona la opción "\textit{Eliminar}" (aparecerá una ventana emergente para confirmar la operación) y pulsa en "\textit{Confirmar}".
		\end{enumerate}\\
		\textbf{Postcondición}        &  - \\
		\textbf{Excepciones}          &  El usuario administrador, no puede modificar o eliminar la cuenta de otro usuario administrador. \\
		\textbf{Importancia}          & Alta \\
		\bottomrule
	\end{tabularx}
	\caption{CU-4 Gestión de usuarios.}
\end{table}

\begin{table}[p]
	\centering
	\begin{tabularx}{\linewidth}{ p{0.21\columnwidth} p{0.71\columnwidth} }
		\toprule
		\textbf{CU-5}    & \textbf{Menú de la aplicación}\\
		\toprule
		\textbf{Requisitos asociados} & RF5 \\
		\textbf{Descripción}          & Permite al usuario navegar por las distintas funcionalidades que tiene la aplicación. \\
		\textbf{Precondición}         & Para poder navegar por todas las funcionalidades es necesario visualizar la aplicación desde una pantalla mediana/grande. \\
		\textbf{Acciones}             &
		\begin{enumerate}
			\def\labelenumi{\arabic{enumi}.}
			\tightlist
			\item El usuario puede acceder al main pulsando "\textit{StatsGlowMind}".
			\item El usuario puede acceder a las funcionalidades que contine el apartado NBA pulsando "\textit{NBA}".
            \item El usuario puede acceder a las funcionalidades que contine el apartado Equipos pulsando "\textit{Equipos}".
            \item El usuario puede acceder a las funcionalidades que contine el apartado Análisis pulsando "\textit{Análisis}".
            \item El usuario puede acceder a las funcionalidades que contine el apartado Predicción "\textit{Predicción}".
            \item El usuario puede acceder al inicio de sesión pulsando "\textit{Iniciar Sesión}".
            \item El usuario puede cerrar sesión pulsando "\textit{Cerrar Sesión}".
		\end{enumerate}\\
		\textbf{Postcondición}        &  - \\
		\textbf{Excepciones}          &  Visualizar la aplicación en versión móvil, en este caso solo hay acceso al main, inicio de sesión y cerrar sesión. \\
		\textbf{Importancia}          & Media \\
		\bottomrule
	\end{tabularx}
	\caption{CU-5 Menú de la aplicación.}
\end{table}

\begin{table}[p]
	\centering
	\begin{tabularx}{\linewidth}{ p{0.21\columnwidth} p{0.71\columnwidth} }
		\toprule
		\textbf{CU-6}    & \textbf{Main de la aplicación}\\
		\toprule
		\textbf{Requisitos asociados} & RF6 \\
		\textbf{Descripción}          & Permite al usuario acceder a todas las funcionalidades que tiene la aplicación de manera visual e intuitivo. \\
		\textbf{Precondición}         & El usuario debe estar registrado para acceder a todas las funcionalidades. \\
		\textbf{Acciones}             &
		\begin{enumerate}
			\def\labelenumi{\arabic{enumi}.}
			\tightlist
			\item El usuario puede acceder a las funcionalidades del apartado NBA pulsando "\textit{NBA}".
            \item El usuario puede acceder a las funcionalidades del apartado Equipos pulsando "\textit{Equipos}".
            \item El usuario puede acceder a las funcionalidades del apartado Análisis pulsando "\textit{Análisis}".
            \item El usuario puede acceder a las funcionalidades del apartado Predicción "\textit{Predicción}".
		\end{enumerate}\\
		\textbf{Postcondición}        &  - \\
		\textbf{Excepciones}          &  Un usuario invitado no puede acceder a todas las funcionalidades. \\
		\textbf{Importancia}          & Media \\
		\bottomrule
	\end{tabularx}
	\caption{CU-6 Main de la aplicación.}
\end{table}

\begin{table}[p]
	\centering
	\begin{tabularx}{\linewidth}{ p{0.21\columnwidth} p{0.71\columnwidth} }
		\toprule
		\textbf{CU-7}    & \textbf{Footer de la aplicación}\\
		\toprule
		\textbf{Requisitos asociados} & RF7, RF7.1, RF7.2, RF7.3, RF7.4 \\
		\textbf{Descripción}          & Permite al usuario acceder a información acerca de la aplicación. \\
		\textbf{Precondición}         & - \\
		\textbf{Acciones}             &
		\begin{enumerate}
			\def\labelenumi{\arabic{enumi}.}
			\tightlist
			\item El usuario puede acceder a la información del desarrollo de la aplicación.
            \item El usuario puede acceder a la información del diseño de la aplicación.
            \item El usuario puede acceder a la información acerca de la aplicación.
            \item El usuario puede acceder a las redes sociales de la aplicación.
		\end{enumerate}\\
		\textbf{Postcondición}        &  - \\
		\textbf{Excepciones}          &  - \\
		\textbf{Importancia}          & Baja \\
		\bottomrule
	\end{tabularx}
	\caption{CU-7 Footer de la aplicación.}
\end{table}

\begin{table}[p]
	\centering
	\begin{tabularx}{\linewidth}{ p{0.21\columnwidth} p{0.71\columnwidth} }
		\toprule
		\textbf{CU-8}    & \textbf{Mostrar información en tiempo real}\\
		\toprule
		\textbf{Requisitos asociados} & RF8, RF8.1, RF8.2, RF8.3, RF8.4 \\
		\textbf{Descripción}          & Permite al usuario visualizar información y estadísticas detalladas sobre la actual temporada de la NBA, en tiempo real. \\
		\textbf{Precondición}         & Tener acceso a la API. \\
		\textbf{Acciones}             &
		\begin{enumerate}
			\def\labelenumi{\arabic{enumi}.}
			\tightlist
            \item El usuario accede a la página "\textit{NBA}".
			\item El usuario puede visualizar la información y estadísticas de la clasificación pulsando en "\textit{Clasificación}".
            \item El usuario puede visualizar la información de los playoffs pulsando en "\textit{Playoffs}".
            \item El usuario puede visualizar la información de los partidos de la temporada pulsando en "\textit{Partidos}".
            \item El usuario puede visualizar la información y estadísticas de los lideres de la liga pulsando en "\textit{Líderes}".
		\end{enumerate}\\
		\textbf{Postcondición}        &  - \\
		\textbf{Excepciones}          &  Se hayan consumido todas las peticiones a la API. \\
		\textbf{Importancia}          &  Alta \\
		\bottomrule
	\end{tabularx}
	\caption{CU-8 Mostrar información en tiempo real.}
\end{table}

\begin{table}[p]
	\centering
	\begin{tabularx}{\linewidth}{ p{0.21\columnwidth} p{0.71\columnwidth} }
		\toprule
		\textbf{CU-9}    & \textbf{Mostrar partidos}\\
		\toprule
		\textbf{Requisitos asociados} & RF8.3 \\
		\textbf{Descripción}          & Permite al usuario visualizar los partidos que se van a disputar y los que se han disputado, junto con sus resultados. \\
		\textbf{Precondición}         & Tener acceso a la API. \\
		\textbf{Acciones}             &
		\begin{enumerate}
			\def\labelenumi{\arabic{enumi}.}
			\tightlist
			\item El usuario accede a los partidos pulsando en "\textit{Partidos}".
                \item El usuario visualiza los partidos jugados o por jugar de la fecha actual.
                \item El usuario puede visualizar partidos jugados o por jugar de la fecha deseada, seleccionado la fecha en el calendario.
		\end{enumerate}\\
		\textbf{Postcondición}        &  - \\
		\textbf{Excepciones}          &  Se hayan consumido todas las peticiones a la API. \\
		\textbf{Importancia}          &  Alta \\
		\bottomrule
	\end{tabularx}
	\caption{CU-9 Mostrar partidos.}
\end{table}

\begin{table}[p]
	\centering
	\begin{tabularx}{\linewidth}{ p{0.21\columnwidth} p{0.71\columnwidth} }
		\toprule
		\textbf{CU-10}    & \textbf{Mostrar líderes}\\
		\toprule
		\textbf{Requisitos asociados} & RF8.4 \\
		\textbf{Descripción}          & Permite al usuario visualizar los líderes de la NBA junto con sus estadísticas. \\
		\textbf{Precondición}         & Tener acceso a la API. \\
		\textbf{Acciones}             &
		\begin{enumerate}
			\def\labelenumi{\arabic{enumi}.}
			\tightlist
			\item El usuario accede a los líderes pulsando en "\textit{Líderes}".
            \item El usuario visualiza los líderes en puntos por partido de la NBA.
            \item El usuario puede visualizar líderes de otra categoría seleccionando la categoría deseada.
		\end{enumerate}\\
		\textbf{Postcondición}        &  - \\
		\textbf{Excepciones}          &  Se hayan consumido todas las peticiones a la API. \\
		\textbf{Importancia}          &  Alta \\
		\bottomrule
	\end{tabularx}
	\caption{CU-10 Mostrar líderes.}
\end{table}

\begin{table}[p]
	\centering
	\begin{tabularx}{\linewidth}{ p{0.21\columnwidth} p{0.71\columnwidth} }
		\toprule
		\textbf{CU-11}    & \textbf{Visualización de equipos}\\
		\toprule
		\textbf{Requisitos asociados} & RF9 \\
		\textbf{Descripción}          & Permite al usuario visualizar todos los equipos de la NBA \\
		\textbf{Precondición}         & El usuario debe estar registrado. \\
		\textbf{Acciones}             & El usuario accede a la página "\textit{Equipos}". \\
		\textbf{Postcondición}        &  Tener acceso a la base de datos. \\
		\textbf{Excepciones}          &  No cargue la base de datos. \\
		\textbf{Importancia}          &  Alta \\
		\bottomrule
	\end{tabularx}
	\caption{CU-11 Visualización de equipos.}
\end{table}

\begin{table}[p]
	\centering
	\begin{tabularx}{\linewidth}{ p{0.21\columnwidth} p{0.71\columnwidth} }
		\toprule
		\textbf{CU-12}    & \textbf{Estadísticas de equipos}\\
		\toprule
		\textbf{Requisitos asociados} & RF9.1 \\
		\textbf{Descripción}          & Permite al usuario visualizar mediante tablas, todas las estadísticas disponibles del equipo seleccionado. \\
		\textbf{Precondición}         & El usuario debe estar registrado. \\
		\textbf{Acciones}             & El usuario selecciona "\textit{Estadísticas}" del equipo deseado. \\
		\textbf{Postcondición}        &  Tener acceso a la base de datos. \\
		\textbf{Excepciones}          &  No cargue la base de datos. \\
		\textbf{Importancia}          &  Alta \\
		\bottomrule
	\end{tabularx}
	\caption{CU-12 Estadísticas de equipos.}
\end{table}

\begin{table}[p]
	\centering
	\begin{tabularx}{\linewidth}{ p{0.21\columnwidth} p{0.71\columnwidth} }
		\toprule
		\textbf{CU-13}    & \textbf{Jugadores del equipo}\\
		\toprule
		\textbf{Requisitos asociados} & RF9.2 \\
		\textbf{Descripción}          & Permite al usuario visualizar mediante una tabla, información general sobre todos los jugadores del equipo seleccionado. \\
		\textbf{Precondición}         & El usuario debe estar registrado. \\
		\textbf{Acciones}             & El usuario selecciona "\textit{Jugadores}" del equipo deseado. \\
		\textbf{Postcondición}        &  Tener acceso a la base de datos. \\
		\textbf{Excepciones}          &  No cargue la base de datos. \\
		\textbf{Importancia}          &  Alta \\
		\bottomrule
	\end{tabularx}
	\caption{CU-13 Jugadores del equipo.}
\end{table}

\begin{table}[p]
	\centering
	\begin{tabularx}{\linewidth}{ p{0.21\columnwidth} p{0.71\columnwidth} }
		\toprule
		\textbf{CU-14}    & \textbf{Mostrar jugadores}\\
		\toprule
		\textbf{Requisitos asociados} & RF10, R10.1 \\
		\textbf{Descripción}          & Permite al usuario visualizar mediante una tabla, todos los jugadores de la NBA y mediante cards los jugadores más destacados de la temporada. \\
		\textbf{Precondición}         & El usuario debe estar registrado. \\
		\textbf{Acciones}             & El usuario accede a la página "\textit{Análisis}". \\
		\textbf{Postcondición}        &  Tener acceso a la base de datos. \\
		\textbf{Excepciones}          &  No cargue la base de datos. \\
		\textbf{Importancia}          &  Alta \\
		\bottomrule
	\end{tabularx}
	\caption{CU-14 Mostrar jugadores.}
\end{table}

\begin{table}[p]
	\centering
	\begin{tabularx}{\linewidth}{ p{0.21\columnwidth} p{0.71\columnwidth} }
		\toprule
		\textbf{CU-15}    & \textbf{Gráficos y estadísticas del jugador}\\
		\toprule
		\textbf{Requisitos asociados} & RF10, R10.2 \\
		\textbf{Descripción}          & Permite al usuario visualizar mediante una tabla, las estadísticas medias de la temporada actual y analizar mediante gráficas de varios tipos las estadísticas del jugador. \\
		\textbf{Precondición}         &  El usuario debe estar registrado. \\
		\textbf{Acciones}             &  El usuario selecciona el jugador deseado. \\
		\textbf{Postcondición}        &  
            \begin{enumerate}
			\def\labelenumi{\arabic{enumi}.}
			\tightlist
			\item Tener acceso a la base de datos.
                \item El jugador seleccionado debe haber jugado al menos un partido de la temporada.
		\end{enumerate}\\
		\textbf{Excepciones}          &  No cargue la base de datos. \\
		\textbf{Importancia}          &  Alta \\
		\bottomrule
	\end{tabularx}
	\caption{CU-15 Gráficos y estadísticas del jugador.}
\end{table}

\begin{table}[p]
	\centering
	\begin{tabularx}{\linewidth}{ p{0.21\columnwidth} p{0.71\columnwidth} }
		\toprule
		\textbf{CU-16}    & \textbf{Predicción de partidos}\\
		\toprule
		\textbf{Requisitos asociados} & RF11 \\
		\textbf{Descripción}          & Permite al usuario seleccionar dos equipos para enfrentarse en un futuro partido y recibir una predicción de dicho partido. \\
		\textbf{Precondición}         &  El usuario debe estar registrado.\\
		\textbf{Acciones}             &  
            \begin{enumerate}
			\def\labelenumi{\arabic{enumi}.}
			\tightlist
			\item El usuario pulsa en "\textit{Predicciones}".
                \item El usuario selecciona los equipos deseados para el partido.
                \item El usuario pulsa el botón "\textit{Predecir Partido}".
		\end{enumerate}\\
		\textbf{Postcondición}        &  Se deben seleccionar los dos equipos. \\
		\textbf{Excepciones}          &  No cargue la función de Cloud Functions. \\
		\textbf{Importancia}          &  Alta \\
		\bottomrule
	\end{tabularx}
	\caption{CU-16 Predicción de partidos.}
\end{table}

\begin{table}[p]
	\centering
	\begin{tabularx}{\linewidth}{ p{0.21\columnwidth} p{0.71\columnwidth} }
		\toprule
		\textbf{CU-17}    & \textbf{Cambio de idioma}\\
		\toprule
		\textbf{Requisitos asociados} & RF12 \\
		\textbf{Descripción}          & Permite al usuario cambiar el idioma a español o inglés.\\
		\textbf{Precondición}         &  - \\
		\textbf{Acciones}             &  El usuario selecciona el idioma deseado desde un selector (ubicado en el pie de página). \\
		\textbf{Postcondición}        &  - \\
		\textbf{Excepciones}          &  - \\
		\textbf{Importancia}          &  Alta \\
		\bottomrule
	\end{tabularx}
	\caption{CU-17 Cambio de idioma.}
\end{table}

\clearpage

\subsection{Diagrama de casos de uso}

\imagen{anexos/diagrams/diagrama_casos_uso}{Diagrama de Casos de Uso}{1}
\apendice{Especificación de diseño}

\section{Introducción}
En este apartado se van a definir las distintas estructuras de diseño que se han empleado en el proyecto.

\hfill

\section{Diseño de datos}
En la base de datos se pueden observar las siguientes colecciones:
\begin{itemize}
\tightlist
    \item \textbf{Users:} Esta colección contiene los datos de los usuarios registrados en la aplicación, tan solo los datos: id, nombre, email y rol del usuario.
    \item \textbf{Games:} En esta colección se almacenan las estadísticas de cada partido jugado de la temporada 23/24 de la NBA, incluyendo la id del partido y la id de los equipos que se enfrentaron.
    \item \textbf{Teams:} Guarda información genérica y estadísticas media de la temporada de cada uno de los equipos de la NBA.
    \item \textbf{Players:} Esta colección guarda los datos personales y las estadísticas de cada partido jugado en la temporada, de cada jugador de la NBA.
\end{itemize}
Ya que es una base de datos NoSQL, estas colecciones no están relacionadas unas con otras.

\clearpage

\section{Diseño procedimental}

\hfill

\subsection{Registro y Login}

\hfill

Diagrama de secuencia que muestra los pasos que realiza el programa cuando un usuario se registra en la aplicación.
\imagen{anexos/diagrams/secuencia_register_login}{Diagrama de secuencia - Registro y Login}{1}

\clearpage

\subsection{Mostrar información en tiempo real}

\hfill

Este diagrama de secuencia visualiza el funcionamiento de la aplicación al mostrar información de la NBA en tiempo real.
\imagen{anexos/diagrams/secuencia_nba}{Diagrama de secuencia - Mostrar información en tiempo real}{1}

\clearpage

\subsection{Mostrar información en tiempo real}

\hfill

Diagrama de secuencia que indica los pasos del funcionamiento de la aplicación al mostrar equipos, junto con sus estadísticas y jugadores. Incluyendo además el funcionamiento del programa al acceder por primera vez en estas páginas.
\imagen{anexos/diagrams/secuencia_equipos}{Diagrama de secuencia - Mostrar equipos}{1}

\clearpage

\subsection{Análisis de jugadores}

\hfill

En este diagrama de secuencia se pueden observar los pasos que realiza la aplicación al analizar jugadores.
\imagen{anexos/diagrams/secuencia_analisis}{Diagrama de secuencia - Análisis de jugadores}{1}

\vfill

\subsection{Predicción de partidos}

\hfill

Este primer diagrama de secuencia indica el proceso que realiza la aplicación para recoger datos de la base de datos, exportarlos a CSV y almacenarlos en Cloud Storage de Firebase.
\imagen{anexos/diagrams/secuencia_export_csv}{Diagrama de secuencia - Exportar CSV}{1}

\hfill

En este segundo diagrama de secuencia se observa los pasos que realiza la aplicación al realizar la funcionalidad de predicción de partidos.
\imagen{anexos/diagrams/secuencia_predicciones}{Diagrama de secuencia - Predicción de partidos}{1}

\clearpage

\section{Diseño arquitectónico}

Ya que en mi proyecto he utilizado Angular como framework, sigo un diseño arquitectónico basado en componentes junto con el enfoque MVVM (Model-View-ViewModel) que proporciona una estructura clara y modular para desarrollar aplicaciones web.

La arquitectura de Angular \cite{patron} se basa en componentes organizados en módulos. Los bloques de construcción básicos son los componentes, servicios y módulos, que están organizados en conjuntos funcionales llamados NgModules.

Los conceptos clave de esta arquitectura son:
\begin{itemize}
\tightlist
    \item \textbf{Módulos, Componentes y Servicios:} Son los bloques de construcción fundamentales de una aplicación Angular.
    \item \textbf{Routing:} Permite la navegación entre diferentes componentes y vistas en la aplicación.
    \item \textbf{Directivas y Pipes:} Proporcionan funcionalidades adicionales para manipular la UI y los datos.
    \item \textbf{Inyección de Dependencias (DI):} Facilita la modularidad y la reutilización del código al permitir la inyección de servicios y dependencias.
\end{itemize}

\hfill

\subsection{Model-View-ViewModel (MVVM)}

El patrón MVVM (Modelo-Vista-Modelo de Vista) \cite{mvvm} es una arquitectura de diseño que se utiliza comúnmente en aplicaciones de interfaz de usuario. Los aspectos más importantes son:

Componentes del patrón MVVM:
\begin{itemize}
\tightlist
    \item \textbf{Modelo:} Representa los datos y la lógica de la aplicación.
    \item \textbf{Vista:} Define la estructura, el diseño y la apariencia de la interfaz de usuario.
    \item \textbf{Modelo de Vista:} Actúa como intermediario entre la vista y el modelo, proporcionando datos y lógica de presentación.
\end{itemize}

\imagen{anexos/patron_mvvm}{Patrón MVVM}{1}

Como se puede observar en la imagen del patrón, las interacciones entre componentes son que la vista conoce al modelo de vista y el modelo de vista conoce al modelo. Y el modelo no conoce al modelo de vista y viceversa, lo que permite que el modelo evolucione independientemente de la vista.

\hfill

Ventajas del patrón MVVM:
\begin{itemize}
\tightlist
    \item \textbf{Adaptabilidad del Modelo:} El modelo de vista sirve como adaptador para el modelo, facilitando cambios en la interfaz de usuario sin modificar el modelo subyacente.
    \item \textbf{Pruebas Unitarias:} Se pueden crear pruebas unitarias para el modelo de vista y el modelo sin depender de la vista.
    \item \textbf{Flexibilidad en el Rediseño:} La interfaz de usuario se puede rediseñar sin afectar al modelo de vista y al modelo.
\end{itemize}

\clearpage

\section{Diseño de interfaces}

En los primeros pasos del proyecto, se realizó un simple prototipo del diseño de las interfaces de la aplicación web. Utilizando la herramienta Figma se pudo representar fácilmente el prediseño de la aplicación.


\imagen{anexos/interfaces/prototipo_main}{Prototipo del main.}{0.6}
\imagen{anexos/interfaces/prototipo_login}{Prototipo del login.}{0.6}
\imagen{anexos/interfaces/prototipo_stats}{Prototipo de stats.}{0.5}
\imagen{anexos/interfaces/prototipo_teams}{Prototipo de teams.}{0.5}
\imagen{anexos/interfaces/prototipo_analysis}{Prototipo de analysis.}{0.5}
\imagen{anexos/interfaces/prototipo_predict}{Prototipo de predict.}{0.6}

\hfill

Haciendo uso de la herramienta Canva, realice un prototipo del logotipo de la aplicación, el cual se veía así:
\imagen{anexos/interfaces/prototipo_logo}{Prototipo del Logotipo.}{0.4}

\vfill

Este prototipo del diseño de las interfaces de la aplicación y el logotipo, fue evolucionando a medida que iba desarrollando el proyecto, cambiando los colores y la estructura, viéndose mucho más moderno e intuitivo.

Para el logotipo me ayudé de la herramienta Microsoft Copilot, que tras una descripción detallada de lo que estaba buscando para el logotipo, me mostró varios logos de los cuales me quedé con uno que posteriormente modifiqué para personalizarlo a mi gusto. Actualmente el logotipo de la aplicación es este:

\imagen{anexos/interfaces/logo}{Logotipo de la Aplicación.}{0.8}

\hfill

El diseño de las interfaces de la aplicación fue mejorando con las implementaciones de Bootstrap, Angular Material y Chart.js. El resultado final de las interfaces es el siguiente:

\imagen{anexos/interfaces/login}{Interfaz Login.}{0.9}
\imagen{anexos/interfaces/main}{Interfaz Main.}{0.85}
\imagen{anexos/interfaces/stats}{Interfaz Stats.}{1}
\imagen{anexos/interfaces/teams}{Interfaz Teams.}{1}
\imagen{anexos/interfaces/analysis}{Interfaz Analysis.}{0.85}
\imagen{anexos/interfaces/predict}{Interfaz Predict.}{1}
\apendice{Documentación técnica de programación}

\section{Introducción}

En este anexo se presenta la documentación técnica de programación del proyecto. Se detallará la estructura de directorios, el manual del programador y los requisitos necesarios para poder ejecutar el proyecto.

\section{Estructura de directorios}

Dentro del repositorio StatsGlowMind de GitHub \cite{github} se pueden observar 3 directorios principales:

\begin{itemize}
\tightlist
    \item
        \textbf{doc:} Directorio donde se encuentra toda la documentación de la memoria y los anexos del proyecto, junto con sus imágenes y bibliografías.
    \item 
        \textbf{predition\_model:} Directorio que contiene el código del desarrollo del modelo Random Forest para la predicción de partidos.
    \item
        \textbf{stats-glow-mind:} Directorio principal donde se halla toda la lógica, el desarrollo y la interfaz del proyecto.
\end{itemize}

A continuación, se va a describir más detalladamente el directorio principal del proyecto \textbf{stats-glow-mind}:
\begin{itemize}
\tightlist
    \item
        \textbf{functions:} Este directorio contiene la función desarrollada en Python \textit{predict} que realiza las predicciones de los partidos de la NBA.
    \item 
        \textbf{src:} Este directorio contiene el código fuente de la aplicación. Todo el desarrollo principal se lleva a cabo aquí.
    \item
        \textbf{src/app:} Contiene el núcleo de la aplicación, incluyendo componentes, servicios, módulos, el enrutamiento de la aplicación y modelos.
    \item
        \textbf{src/app/components:} Contiene todos los componentes reutilizables de la aplicación, cada componente tiene su lógica e interfaz.
    \item
        \textbf{src/app/material:} Contiene un módulo que importa todos los componentes de angular material usados en el proyecto.
    \item
        \textbf{src/app/models:} Contiene distintos modelos de datos utilizados en la aplicación.
    \item
        \textbf{src/app/services:} Contiene todos los servicios que encapsulan la lógica de la aplicación, como las interacciones con las APIs, la base de datos, autenticación de usuarios, \ldots
    \item
        \textbf{src/assets:} Contiene los recursos estáticos de la aplicación como imágenes, temas personalizados e internacionalización del proyecto.
    \item
        \textbf{src/assets/i18n:} Contiene los archivos de traducción de la aplicación, en español y en inglés.
    \item
        \textbf{src/assets/images:} Contiene todas las imágenes, íconos, logos y fondos de pantalla usados en la aplicación.
    \item
        \textbf{src/assets/theme:} Contiene un tema de angular material personalizado con la paleta de colores de la aplicación.
    \item
        \textbf{src/environments:} Contiene archivos de configuración de entorno para diferentes entornos (producción y desarrollo).
\end{itemize}

\hfill

\section{Manual del programador}

Este manual proporciona información detallada para los desarrolladores que deseen contribuir al proyecto, como el proceso de instalación, configuración, ejecución y despliegue.

\subsection{Instalación de herramientas}

Antes de empezar a montar el proyecto en local, es imprescindible contar con 3 herramientas que son necesarias instalar.

\subsubsection{Visual Studio Code}
Visual Studio Code (VS Code) es un editor de código fuente, es el que yo he utilizado para desarrollar el proyecto y el que recomiendo usar.

Para instalar Visual Studio Code:
\begin{enumerate}
    \item Descarga Visual Studio Code desde su sitio web oficial: \url{https://code.visualstudio.com/}.
    \item Ejecuta el instalador y sigue las instrucciones en pantalla para completar la instalación.
    \item Una vez instalado, puedes abrir VS Code desde el menú de inicio o desde una terminal escribiendo \textit{code .} en el directorio que se desea editar.
\end{enumerate}

\subsubsection{Node.js y npm}
Node.js es un entorno de ejecución para JavaScript en el lado del servidor, y npm es el gestor de paquetes de Node.js. Aunque no se usa directamente en este proyecto, npm es necesario para gestionar las dependencias de Angular.

Para instalar Node.js y npm:
\begin{enumerate}
    \item Descarga el instalador de Node.js desde su sitio web oficial: \url{https://nodejs.org/en/download/package-manager}.
    \item Ejecuta el instalador y sigue las instrucciones. Este instalador también instalará npm.
    \item Abre una terminal y ejecuta los siguientes comandos para verificar la instalación:
\end{enumerate}
$$ \textit{node -v} $$
$$ \textit{npm -v} $$

\subsubsection{Python}
Python es necesario en este proyecto para crear o modificar funciones de Cloud Functions (el backend del proyecto).

Para instalar Python:
\begin{enumerate}
    \item Descarga el instalador de Python desde su sitio web oficial: \url{https://www.python.org/downloads/}.
    \item Ejecuta el instalador y sigue las instrucciones de la instalación.
    \item Para verificar la instalación, abre una terminal y ejecuta el comando:
\end{enumerate}
$$ \textit{python --version} $$


\subsection{Clonar repositorio GitHub}

El primer paso para poder montar el proyecto en local, es clonar el repositorio del proyecto subido a GitHub. Para ello abre la terminal del ordenador, navega al directorio donde se quiere clonar el repositorio y ejecuta los siguientes comandos:

$$ \textit{git clone https://github.com/roberlastrass/StatsGlowMind.git} $$
$$ \textit{cd StatsGlowMind} $$

\subsection{Instalar las dependencias de Angular}

Dentro del directorio del proyecto, navega al directorio principal \textbf{stats-glow-mind} e instala las dependencias de angular utilizando npm, instalado anteriormente. Los comandos necesarios son los siguientes:

$$ \textit{cd stats-glow-mind} $$
$$ \textit{npm install} $$

Asegúrate de tener la última versión de Angular CLI ejecutando: \textit{npm install -g @angular/cli}

\subsection{Instalar las dependencias de Python}

También es necesario instalar las dependencias de Python, para ello, dentro del directorio principal del proyecto \textbf{stats-glow-mind}, navega al directorio donde se encuentran las funciones backend del proyecto \textbf{functions} e instala las dependencias que se encuentra en \textit{requirements.txt}, asegúrate de estar en un entorno virtual:

$$ \textit{cd functions} $$
$$ \textit{python -m venv env} $$
$$ \textit{.\text{\textbackslash}env\text{\textbackslash}Scripts\text{\textbackslash}activate} $$
$$ \textit{pip install -r requirements.txt} $$

\subsection{Inicializar Firebase}

Para inicializar firebase en un proyecto existente, lo primero es instalar Firebase CLI (si aún no está instalado) con este comando:
$$ \textit{npm install -g firebase-tools} $$

Después inicializar Firebase en el proyecto, iniciando sesión antes, los comandos utilizados para realizar esto son los siguientes, desde el directorio principal del proyecto \textbf{stats-glow-mind}:
$$ \textit{firebase login} $$
$$ \textit{firebase init} $$

Con esto, seleccionas tu proyecto Firebase y lo configuras a tu manera, agregando la configuración en los diferentes entornos de \textit{environments} y añadiendo tus claves y credenciales de Firebase.

\subsection{Despliegue de funciones en la nube}

Cuando se ha creado una función nueva o se ha modificado una existente, el comando para poder subir estos cambios a Cloud Functions es:
$$ \textit{firebase deploy --only functions} $$

\subsection{Ejecución del proyecto en local}

Una vez realizado este proceso, el proyecto se puede ejecutar en local para observar los cambios realizados y comprobar su funcionamiento, con uno de estos dos comandos, desde el directorio principal del proyecto \textbf{stats-glow-mind}:
$$ \textit{ng serve} $$ $$ \textit{npm run start} $$

Con esto, el proyecto se compilará y lo ejecutará localmente en la ruta: \url{http://localhost:4200/}

\subsection{Despliegue de la aplicación}

Finalmente, cuando se ha terminado de configurar y probar el funcionamiento de la aplicación localmente, se puede desplegar la aplicación utilizando Firebase Hosting escribiendo el siguiente comando desde el directorio principal del proyecto \textbf{stats-glow-mind}:
$$ \textit{firebase deploy --only hosting} $$

Esto desplegará la aplicación en Firebase Hosting y te proporcionará una URL donde puedes acceder a tu aplicación en la web. En este caso la URL de la aplicación es: 
\url{https://statsglowmindtfg.web.app/}
\apendice{Documentación de usuario}

\section{Introducción}

En este apéndice se explican los aspectos relacionados con la aplicación, por parte del usuario. Se muestran los requisitos de los usuarios y el manual de uso de la aplicación para el usuario.

\section{Requisitos de usuarios}

La aplicación web \textbf{StatsGlowMind} está desplegada para todo el mundo, lo que hace que no se necesiten muchos requisitos, tan solo es necesario contar con un dispositivo compatible como un ordenador, tablet o móvil. Y que dicho dispositivo cuente con un navegador web que contenga los siguientes requisitos mínimos recomendados de versiones:

\begin{itemize}
\tightlist
    \item Google Chrome: versión 87 o superior.
    \item Mozilla Firefox: versión 84 o superior.
    \item Microsoft Edge: versión 88 o superior.
    \item Safari: versión 14 o superior.
    \item Opera: versión 72 o superior.
\end{itemize}

\section{Instalación}

Como esta aplicación está desplegada para que cualquier usuario la pueda utilizar, no necesitan instalar nada, tan solo deben acceder al enlace \url{https://statsglowmindtfg.web.app/} desde un navegador.

\section{Manual del usuario}

En este manual se va a proporcionar una descripción detallada del uso de la aplicación, con el objetivo de facilitar al usuario a utilizar y comprender el funcionamiento de la aplicación.

Para ello, lo primero es acceder a la aplicación desde un navegador con el enlace mencionado anteriormente \url{https://statsglowmindtfg.web.app/} y comienza el uso de la aplicación.

\subsection{Inicio}

Inicialmente se observa la siguiente pantalla, en ella nos encontramos con una barra de navegación, la página principal (main) de la aplicación y el pie de página (footer) de la aplicación.

\imagen{anexos/interfaces/main}{Página de inicio.}{0.52}

\subsection{Barra de navegación}

La barra de navegación cuenta con el logotipo de la aplicación en el que si lo pulsas puedes acceder a la página principal, un menú desplegable con el que puedes acceder a las funcionalidades principales de la aplicación, un botón en el cual se puede iniciar sesión y otro botón para cerrar la sesión iniciada.

\imagen{anexos/user_manual/barra_navegacion}{Barra de navegación.}{1}

Esta barra va a estar presente en todas las páginas de la aplicación, por lo que podemos acceder a las funcionalidades mencionadas anteriormente en cualquier momento, excepto si utilizamos la aplicación en una pantalla pequeña como puede ser la de un móvil, en la cual la barra de navegación se vería así y tan solo podríamos acceder a la página principal, al inicio de sesión y a cerrar sesión.

\imagen{anexos/user_manual/barra_navegacion_movil}{Barra de navegación móvil.}{0.8}

Los botones de \textit{Iniciar Sesión} y \textit{Cerrar Sesión}, no tienen confusión, ya que si pasas el cursor por encima del botón te muestra el nombre de la función que realiza.

\imagen{anexos/user_manual/tooltip}{Información botones.}{0.3}

\subsection{Página principal}

En la página principal nos encontramos con dos Tabs, que son unos contenedores que contienen varias pestañas. El primer Tab contiene las pestañas de cada funcionalidad principal como son \textbf{NBA}, \textbf{EQUIPOS}, \textbf{ANÁLISIS} y \textbf{PREDICCIONES}, desde estas pestañas se puede acceder a la funcionalidad indicada. El segundo Tab contiene las pestañas con las funcionalidades de la NBA: \textbf{Clasificación}, \textbf{Playoffs}, \textbf{Partidos} y \textbf{Líderes}.

\imagen{anexos/user_manual/tabs}{Tabs página principal.}{0.9}

Justo debajo de estos Cards, podemos observar varios Cards (tarjetas) con sus descripciones, en los que cada uno de ellos, te lleva a una funcionalidad principal distinta.

\imagen{anexos/user_manual/cards}{Cards página principal.}{1}

También contiene un botón en la parte de arriba a la derecha de la página, en el cual contiene el nombre del usuario que ha iniciado sesión o en caso de que no se haya iniciado sesión, se mostraría como "\textit{Usuario Invitado}". Si pulsas en este botón y eres un usuario con permisos de administrador, accedes a la gestión de usuarios.

\imagen{anexos/user_manual/nombre_usuario}{Nombre de usuario.}{0.2}

\hfill

\subsection{Pie de página}

El pie de página se muestra en todas las páginas de la aplicación y en él se puede observar información sobre el \textbf{DESARROLLO}, \textbf{DISEÑO}, \textbf{ACERCA DE} y \textbf{REDES} de la aplicación. Cada información de cada uno de estos apartados te lleva a una página distinta mediante un enlace.

\imagen{anexos/user_manual/footer}{Pie de página.}{1}

Abajo de todo a la derecha se encuentra el selector de idioma, en el que puedes cambiar el idioma de español a inglés y viceversa.

\imagen{anexos/user_manual/idioma}{Selector de idioma.}{0.3}

Una vez seleccionas el idioma deseado, se cambia todo el contenido de la aplicación al idioma seleccionado. Así se ve el pie de página cuando seleccionas \textit{Inglés} como idioma de la aplicación.

\imagen{anexos/user_manual/footer_ingles}{Pie de página en inglés.}{1}

\hfill

\subsection{Inicio de Sesión}

Para navegar al inicio de sesión, se accede desde la barra de navegación pulsando en el botón de \textit{Iniciar Sesión}, como ya se mencionó. Una vez hemos accedido, podemos ver un formulario en el que se tiene que rellenar el email y la contraseña para poder iniciar sesión, o si se desea acceder con la cuenta de Google (si la tienes) tan solo hay que pulsar en el botón "\textit{Inicia sesión con Google"} y seleccionas tu cuenta de Google.

\imagen{anexos/user_manual/login}{Inicio de sesión.}{0.7}

Si el inicio de sesión ha sido correcto, la aplicación te redirigirá a la página principal (main), para que puedas continuar con el uso de la aplicación.

Si, por el contrario, el inicio no ha sido válido porque no se ha introducido bien algún dato o el usuario no está registrado aún, te saldrá el siguiente mensaje de error:
\imagen{anexos/user_manual/error_login}{Error de inicio de sesión.}{0.3}

En caso de que el usuario no se haya registrado todavía, tiene la opción de registrarse pulsando en el enlace "\textit{Regístrate}", que se encuentra en la parte de abajo del formulario.
\imagen{anexos/user_manual/boton_register}{Enlace para registrarse.}{0.3}

\hfill

\subsection{Registro de usuarios}

En el registro de usuarios, nos encontramos con otro formulario en el que hay que rellenar más campos: el nombre de usuario, el email, la contraseña, la contraseña repetida y el rol del usuario en el que a no ser que tengas permisos por el administrados, solo se podrá seleccionar el rol "\textit{Usuario}". También es posible saltarse todo esto si se desea registrarse con la cuenta de Google, pulsando en "\textit{Registrarse con Google"}.

\imagen{anexos/user_manual/register}{Registro de usuarios.}{0.6}

Cuando el usuario se registra exitosamente, se mostrará por pantalla el siguiente mensaje y te redirige a la página principal.
\imagen{anexos/user_manual/exito_register}{Enlace para registrarse.}{0.3}

Sin embargo, si existe algún error al rellenar el formulario, porque falta algún dato o el email está mal escrito o la contraseña es demasiado corta (min 6 caracteres), saldrá el siguiente mensaje de error:
\imagen{anexos/user_manual/error_register}{Error de registro.}{0.3}

Y si por algún casual todos los campos del formulario cumplen con las condiciones de registro, pero las contraseñas son distintas, aparece este mensaje de error:
\imagen{anexos/user_manual/error_password}{Contraseñas distintas.}{0.3}


\subsection{Gestión de usuarios}

Esta funcionalidad solo la pueden realizar los usuarios con rol de administrador. Para acceder a esta página, como comentamos en el apartado de la \textbf{Página Principal}, existe un botón en la parte superior a la derecha en el que aparece el nombre del usuario, que pulsando dicho botón se accede a la página de gestión de usuarios. En esta página aparecen todos los usuarios con todos sus datos, incluido la uid.

\imagen{anexos/user_manual/admin}{Gestión de usuarios.}{1}

\subsubsection{Modificación de usuarios}

El administrador podrá modificar cualquier usuario clicando el botón "\textit{MODIFICAR}", de tal forma que se mostrará una pantalla emergente en la que puedes cambiar el nombre y el email del usuario. Y pulsando en "\textit{Confirmar}" se guardan los cambios en la base de datos. En este ejemplo que voy a mostrar se modifica el email del usuario "\textit{lastras}":

\imagen{anexos/user_manual/user_update}{Modificación de usuarios.}{0.7}

Resultado de la modificación:
\imagen{anexos/user_manual/updated_user}{Usuario modificado.}{0.4}

\subsubsection{Eliminación de usuarios}

El administrador también podrá eliminar usuarios clicando el botón "\textit{ELIMINAR}", seguidamente aparece una pantalla emergente para confirmar si de verdad se desea eliminar al usuario, en la que se tendrá que pulsar en "\textit{Confirmar}" para eliminar al usuario seleccionado. En este caso se va a eliminar al usuario "\textit{lastras}":

\imagen{anexos/user_manual/user_delete}{Eliminación de usuarios.}{0.5}

Resultado de la modificación:
\imagen{anexos/user_manual/deleted_user}{Usuario eliminado.}{0.9}


\subsection{NBA}

Esta página que muestra estadísticas a tiempo real de la NBA está compuesta por 4 páginas más: \textbf{Clasificación}, \textbf{Playoffs}, \textbf{Partidos} y \textbf{Líderes}. Para acceder a cada una de estas funcionalidades, se puede realizar desde la barra de navegación (en el menú desplegable) o desde la página principal (en el Tab de la NBA o directamente en el Card NBA).

Se puede acceder a las distintas páginas de la NBA, desde este menú:
\imagen{anexos/user_manual/stats_menu}{Menú NBA.}{1}

\subsubsection{Clasificación}

Por defecto, si accedes a \textit{NBA} desde el Card de la página principal, navegas directamente a esta página, donde se muestra la clasificación de las dos conferencias de la NBA.

\imagen{anexos/user_manual/standings}{Clasificación de la NBA.}{0.72}

El icono naranja de las tablas, es un icono de información que te lleva al glosario de la aplicación para poder entender los términos de baloncesto empleados.
\imagen{anexos/user_manual/glossary_icon}{Icono del glosario.}{0.2}


\subsubsection{Playoffs}

En esta página se observa la tabla eliminatoria de los playoffs de la NBA.

\imagen{anexos/user_manual/playoffs}{Playoffs de la NBA.}{1}


\subsubsection{Partidos}

Desde esta página se pueden observar los próximos partidos que se van a disputar, como los partidos ya finalizados con sus resultados.

\imagen{anexos/user_manual/games}{Partidos de la NBA.}{1}

Para visualizar partidos de otras fechas, hay un calendario en el que se puede seleccionar el día deseado y observar los partidos jugados ese día. Todos estos partidos se muestran en la zona horaria UTC.

\imagen{anexos/user_manual/calendar}{Calendario de partidos.}{0.6}

Si en la fecha seleccionada no se jugó ningún partido, entonces aparecerá el siguiente texto:
\imagen{anexos/user_manual/no_games}{No se encontraron partidos.}{0.7}

\hfill

\subsubsection{Líderes}

Esta página muestra mediante una tabla las estadísticas de los jugadores líderes de la NBA según la categoría: Puntos, Rebotes, Asistencias, Robos, Tapones o Eficiencia.

\imagen{anexos/user_manual/leaders}{Líderes de la NBA.}{1}

Para cambiar de categoría, existe un selector en el que puedes seleccionar la categoría deseada:
\imagen{anexos/user_manual/category}{Categorías.}{0.5}

Al final de la tabla nos encontramos un paginador en el que podemos seleccionar cuantas filas queremos que se muestren y donde podemos cambiar de página:
\imagen{anexos/user_manual/paginator}{Paginador de la tabla.}{0.6}

\clearpage

\subsection{Equipos}

En la página de \textit{Equipos} se visualizan todos los equipos de la NBA y se puede acceder a las estadísticas de cada equipo pulsando en "\textit{ESTADÍSTICAS"} y a los jugadores de cada equipo pulsando en \textit{"JUGADORES"} del equipo deseado.

\imagen{anexos/user_manual/teams}{Equipos de la NBA.}{1}

\subsubsection{Estadísticas del equipo}

Desde esta página vemos las estadísticas del equipo incluyendo las estadísticas totales, en casa, en visitante, en partidos ganados y en partidos perdidos. También es posible acceder a los jugadores del mismo equipo pulsando al botón "\textit{JUGADORES"} ubicado en la parte superior de la derecha.
\imagen{anexos/user_manual/team_stats}{Estadísticas del equipo.}{1}

\subsubsection{Jugadores del equipo}

En esta página vemos los datos generales de los jugadores del equipo. Es posible acceder a las estadísticas del mismo equipo pulsando al botón "\textit{ESTADÍSTICAS"} ubicado en la parte superior de la izquierda.
\imagen{anexos/user_manual/team_players}{Jugadores del equipo.}{0.9}

\subsection{Análisis}

Cuando navegamos a la página de Análisis de los jugadores de la NBA, lo primero que nos aparece son unas Cards con los jugadores más destacados de la temporada. Justo debajo encontramos una tabla en la que se encuentran todos los jugadores de la NBA ordenados alfabéticamente.

\imagen{anexos/user_manual/analysis}{Análisis de jugadores.}{0.7}

En la parte superior de la tabla se encuentra un buscador, en el que podemos buscar un jugador en específico, como por ejemplo \textit{LeBron James}:
\imagen{anexos/user_manual/search}{Buscador de jugadores.}{1}

Para acceder al análisis del jugador, tan solo hay que clicar el jugador deseado o en caso de las Cards, pulsar en "\textit{ANÁLISIS}".

\subsubsection{Análisis del jugador}

En el análisis del jugador, primero se observa una tabla con las estadísticas totales de la temporada.
\imagen{anexos/user_manual/player_stats}{Estadísticas del jugador.}{1}

Seguidamente tenemos varias gráficas de datos, con los que podemos analizar el rendimiento del jugador.

Primero tenemos una gráfica de líneas en la que muestra las estadísticas principales por partido, como son los puntos, asistencias, rebotes, robos y tapones. En este tipo de gráficas es posible ocultar líneas de estadísticas pulsando directamente en el nombre de la estadística, por ejemplo, en BLK (Tapones).
\imagen{anexos/user_manual/stats_chart}{Gráfica de estadísticas por partido.}{1}

La siguiente gráfica es otra gráfica de líneas en la que se analizan los rebotes por partido, incluyendo rebotes defensivos, ofensivos y totales.
\imagen{anexos/user_manual/rebounds_chart}{Gráfica de rebotes por partido.}{1}

Continuamos con unas gráficas circulares, en las que se analizan los tiros de campo del jugador. El primer gráfico circular muestra los porcentajes de los puntos de media según el tipo de tiro de campo: tiro de dos puntos, triples y tiros libres. El resto de gráficos muestran el porcentaje de tiros de campo anotados y fallados, según el tipo de tiro.
\imagen{anexos/user_manual/field_goals_chart}{Gráfica de tiros de campo.}{0.95}

Por último, se muestra una gráfica de barras que indica la diferencia de puntos cuando el jugador está en la pista.
\imagen{anexos/user_manual/plus_minus_chart}{Gráfica de +/-.}{1}


\subsection{Predicciones}

Desde esta página se realizan las predicciones, para ello es necesario seleccionar los dos equipos que se van a enfrentar y después pulsar en "\textit{Predecir Partido}".

\imagen{anexos/user_manual/predict}{Predicciones de partidos.}{0.9}

Tras varios segundos esperando a que se calcule la predicción del partido, se mostrará en la parte inferior una gráfica circular en la que se muestran las probabilidades de victoria de cada equipo y al lado el equipo ganador según la predicción.

\imagen{anexos/user_manual/winner}{Equipo ganador.}{0.9}

En caso de que no se hayan seleccionado los dos equipos que se van a enfrentar, saldrá el siguiente mensaje de error:
\imagen{anexos/user_manual/error_predict}{Error de predicciones.}{0.5}

\subsection{Acerca de StatsGlowMind}

Esta página proporciona información acerca de la aplicación, así como, \textit{¿Qué es StatsGlowMind?}, \textit{Desarrollo}, \textit{Diseño} y \textit{Redes Sociales} de la aplicación.
\imagen{anexos/user_manual/about}{Acerca de StatsGlowMind.}{0.9}

\subsection{Glosario}
En el \textit{Glosario} de la aplicación se visualiza una tabla con los términos empleados en estadísticas y gráficos de baloncesto, de esta forma, el usuario que no conozca algún término, podrá entender toda la información de la aplicación sin problema.
\imagen{anexos/user_manual/glossary}{Glosario de la aplicación.}{0.7}

\apendice{Anexo de sostenibilización curricular}

\section{Introducción}
Este anexo incluye una reflexión personal sobre los aspectos de sostenibilidad abordados en el desarrollo de este trabajo. El objetivo es mostrar cómo se han adquirido y aplicado competencias de sostenibilidad durante el Trabajo de Fin de Grado, conforme a las directrices de la CRUE. A continuación, se detallan las competencias adquiridas y su aplicación práctica en el proyecto.

\section{Competencias de sostenibilidad}

\subsection{Conciencia y comprensión de la sostenibilidad}
Durante el desarrollo del proyecto, he adquirido una mayor conciencia sobre la importancia de la sostenibilidad en el ámbito tecnológico. Entender que las decisiones técnicas pueden tener un impacto significativo en el medio ambiente y la sociedad es crucial. Por ejemplo, la elección de servicios de almacenamiento y procesamiento en la nube fue influenciada por la búsqueda de proveedores con políticas de energía renovable y eficiencia energética. Firebase, como plataforma de backend, proporciona una infraestructura eficiente y escalable, minimizando el desperdicio de recursos.

\subsection{Diseño sostenible}
El diseño de la aplicación web de la NBA se ha realizado teniendo en cuenta principios de sostenibilidad. Se ha optimizado el código para reducir el consumo de recursos y mejorar la eficiencia. La implementación de algoritmos de Machine Learning, como el modelo Random Forest, se ha llevado a cabo de manera que el procesamiento se realice de forma eficiente, reduciendo la carga computacional y, por ende, el consumo energético.

\subsection{Uso responsable de los recursos}
La aplicación utiliza APIs externas para obtener datos en tiempo real. Para minimizar el número de peticiones y, por lo tanto, el consumo de recursos, se han implementado mecanismos de caché y almacenamiento temporal de datos. Esto no solo mejora la eficiencia de la aplicación, sino que también reduce la carga en los servidores externos, contribuyendo a un uso más responsable y sostenible de los recursos disponibles.

\subsection{Impacto social y accesibilidad}
La sostenibilidad también implica considerar el impacto social de la tecnología. En este proyecto, se ha trabajado para asegurar que la aplicación sea accesible para un público diverso. La implementación de funcionalidades de internacionalización (i18n) y la inclusión de traducciones con ngx-translate garantizan que usuarios de diferentes idiomas y regiones puedan utilizar la aplicación sin barreras lingüísticas. Además, se ha prestado atención a la accesibilidad web, asegurando que la aplicación cumpla con los estándares de accesibilidad y sea usable para personas con discapacidades.

\subsection{Educación y sensibilización}
El desarrollo del proyecto ha sido una oportunidad para educar y sensibilizar sobre la sostenibilidad en el ámbito tecnológico. Compartir las experiencias y los desafíos encontrados durante el desarrollo del proyecto con otros estudiantes y profesionales ha contribuido a crear una mayor conciencia sobre la importancia de la sostenibilidad en la tecnología. La documentación y el código del proyecto se han hecho disponibles de manera abierta para fomentar la colaboración y el aprendizaje colectivo.

\subsection{Evaluación y mejora continua}
Finalmente, la sostenibilidad es un proceso continuo que requiere evaluación y mejora constante. Durante el desarrollo del proyecto, se ha establecido un ciclo de retroalimentación para evaluar el impacto de las decisiones tomadas y buscar continuamente formas de mejorar. Esto incluye la monitorización del rendimiento de la aplicación y la búsqueda de nuevas tecnologías y prácticas que puedan hacer el proyecto más sostenible.

\section{Conclusión}
La integración de principios de sostenibilidad en el desarrollo de este proyecto ha sido profundamente educativa, permitiéndome valorar la importancia de considerar el impacto ambiental y social de las decisiones técnicas. A través de la optimización del uso de recursos, la elección de proveedores con políticas de energía renovable y la atención a la accesibilidad e impacto social, he aprendido que es posible crear soluciones tecnológicas avanzadas y responsables. Este enfoque no solo beneficia al medio ambiente y a la sociedad, sino que también añade valor al proyecto, haciéndolo más robusto y alineado con las expectativas actuales de sostenibilidad. Las competencias adquiridas durante este trabajo no solo son valiosas para el éxito de este proyecto, sino que también representan una base sólida para futuros proyectos profesionales. La sostenibilidad en la tecnología no es solo una tendencia, sino una necesidad para construir un futuro mejor y más equitativo.


\bibliographystyle{plain}
\bibliography{bibliografiaAnexos}

\end{document}
